%% Working draft. Structure is generated; prose is authored here.
%% Build from the repository root with: bash scripts/build_paper.sh <this paper dir>
\documentclass[11pt]{article}

\usepackage[utf8]{inputenc}
\usepackage[T1]{fontenc}
\usepackage{amsmath,amssymb}
\usepackage{booktabs}
\usepackage{graphicx}
\usepackage{tikz}
\usepackage[margin=1in]{geometry}
\usepackage[hidelinks]{hyperref}

\graphicspath{{../figs/out/}}

\renewcommand{\topfraction}{0.9}
\renewcommand{\bottomfraction}{0.7}
\renewcommand{\textfraction}{0.08}
\renewcommand{\floatpagefraction}{0.8}
\setcounter{topnumber}{2}
\setcounter{totalnumber}{3}

%% Numbers and tables are emitted by the analysis code from hash-bound evidence.
% Generated by the figure script from hash-bound evidence. Do not hand-edit.
\newcommand{\AlphaLevel}{0.05}
\newcommand{\Questions}{38}
\newcommand{\Conditions}{3}
\newcommand{\ConditionsWord}{three}
\newcommand{\Graders}{three}
\newcommand{\Answers}{228}
\newcommand{\AnswersPerSet}{114}
\newcommand{\SubsetN}{18}
\newcommand{\SubsetNWord}{eighteen}
\newcommand{\TextOnlyN}{20}
\newcommand{\Contrasts}{18}
\newcommand{\Model}{qwen2.5:3b}
\newcommand{\Temperature}{0}
\newcommand{\HarmMM}{0.389}
\newcommand{\HarmText}{0.444}
\newcommand{\HarmDiff}{$-5.6$}
\newcommand{\HarmP}{1.000}
\newcommand{\HarmDiscordant}{7}
\newcommand{\HarmK}{one}
\newcommand{\HarmCandidates}{19}
\newcommand{\HarmGradings}{36}
\newcommand{\HarmLooSign}{14}
\newcommand{\HarmLooBreaks}{4}
\newcommand{\HelpMM}{0.556}
\newcommand{\HelpText}{0.222}
\newcommand{\HelpDiff}{$+33.3$}
\newcommand{\HelpP}{0.070}
\newcommand{\HelpDiscordant}{8}
\newcommand{\HelpK}{six}
\newcommand{\HelpVerdictK}{one}
\newcommand{\HelpVerdictCandidates}{12}
\newcommand{\HelpFloor}{0.008}
\newcommand{\HelpLooVerdict}{17}
\newcommand{\SigDiff}{$+60.0$}
\newcommand{\SigP}{$4.9 \times 10^{-4}$}
\newcommand{\SigDiscordant}{12}
\newcommand{\SigK}{four}
\newcommand{\SigSignK}{twelve}
\newcommand{\SigHolm}{0.008}
\newcommand{\KillThreshold}{five}
\newcommand{\SmallestK}{one}
\newcommand{\KillState}{did not fire}
\newcommand{\SignificanceNeedsM}{six}
\newcommand{\SweepLow}{0.001}
\newcommand{\SweepHigh}{0.20}
\newcommand{\SweepLevels}{25}
\newcommand{\IncapableSweepMin}{3}
\newcommand{\MedianKSweepMax}{7}
\newcommand{\MedianKSweepMin}{3}
\newcommand{\Incapable}{6}
\newcommand{\AtFloor}{4}
\newcommand{\AtFloorNS}{2}
\newcommand{\AtFloorNSDiscordant}{four}
\newcommand{\SurvivesHolm}{3}
\newcommand{\SignificantRaw}{4}
\newcommand{\Directional}{15}
\newcommand{\FragileMax}{two}
\newcommand{\FragileSign}{4}
\newcommand{\FragileVerdict}{7}
\newcommand{\DecoupledGap}{four}
\newcommand{\DecoupledCount}{9}
\newcommand{\HelpRatio}{six}
\newcommand{\RegradeAnswers}{228}
\newcommand{\RegradeContested}{four}
\newcommand{\RegradeImageFlip}{one}
\newcommand{\RegradeMovingRows}{two}
\newcommand{\RegradeIdenticalRows}{ten}
\newcommand{\ClosedFormMismatches}{0}
\newcommand{\ClosedFormStates}{13,761}
\newcommand{\ClosedFormPairs}{1,201}
\newcommand{\KappaVerInversionsLow}{95}
\newcommand{\KappaVerInversionsHigh}{98}
\newcommand{\KappaVerSigTables}{434}
\newcommand{\SimTables}{840,000}
\newcommand{\SimReps}{10,000}
\newcommand{\FragileGivenSigNEighteen}{0.67}
\newcommand{\FragileGivenSigNThirtyEight}{0.38}
\newcommand{\SpearmanSigNThirtyEight}{$-0.93$}
\newcommand{\NumericHolmLT}{0.066}
\newcommand{\NumericHolmSurvivors}{one}
\newcommand{\LLMKappaQwen}{0.97}
\newcommand{\LLMKappaLlama}{0.91}
\newcommand{\LLMKappaBoth}{0.94}
\newcommand{\RegradedMM}{0.444}
\newcommand{\RegradedDiff}{$+0.0$}
\newcommand{\SummaryHarmStrict}{1}
\newcommand{\SummaryHarmNumeric}{0}
\newcommand{\SummaryHelpStrict}{6}
\newcommand{\SummaryHelpNumeric}{6}
\newcommand{\SummarySigStrict}{4}
\newcommand{\SummarySigNumeric}{2}
\newcommand{\StaticTier}{missing}
\newcommand{\PrimaryTier}{partial}
\newcommand{\SotaTier}{blocked}
\newcommand{\RetrievalN}{38}
\newcommand{\RetrievalCorrect}{38}
\newcommand{\IndexFiles}{0}
\newcommand{\BlockedComponent}{full LLM QA}
\newcommand{\ForgedState}{were not}
\newcommand{\NEvidence}{21}
\newcommand{\EvidenceBytes}{527,252}

%% Draft-only owner placeholder (ruling R15, 2026-09-09). Prints ONLY when the
%% build defines \DRAFTTODO (e.g. pdflatex "\def\DRAFTTODO{1}%% Working draft. Structure is generated; prose is authored here.
%% Build from the repository root with: bash scripts/build_paper.sh <this paper dir>
\documentclass[11pt]{article}

\usepackage[utf8]{inputenc}
\usepackage[T1]{fontenc}
\usepackage{amsmath,amssymb}
\usepackage{booktabs}
\usepackage{graphicx}
\usepackage{tikz}
\usepackage[margin=1in]{geometry}
\usepackage[hidelinks]{hyperref}

\graphicspath{{../figs/out/}}

\renewcommand{\topfraction}{0.9}
\renewcommand{\bottomfraction}{0.7}
\renewcommand{\textfraction}{0.08}
\renewcommand{\floatpagefraction}{0.8}
\setcounter{topnumber}{2}
\setcounter{totalnumber}{3}

%% Numbers and tables are emitted by the analysis code from hash-bound evidence.
% Generated by the figure script from hash-bound evidence. Do not hand-edit.
\newcommand{\AlphaLevel}{0.05}
\newcommand{\Questions}{38}
\newcommand{\Conditions}{3}
\newcommand{\ConditionsWord}{three}
\newcommand{\Graders}{three}
\newcommand{\Answers}{228}
\newcommand{\AnswersPerSet}{114}
\newcommand{\SubsetN}{18}
\newcommand{\SubsetNWord}{eighteen}
\newcommand{\TextOnlyN}{20}
\newcommand{\Contrasts}{18}
\newcommand{\Model}{qwen2.5:3b}
\newcommand{\Temperature}{0}
\newcommand{\HarmMM}{0.389}
\newcommand{\HarmText}{0.444}
\newcommand{\HarmDiff}{$-5.6$}
\newcommand{\HarmP}{1.000}
\newcommand{\HarmDiscordant}{7}
\newcommand{\HarmK}{one}
\newcommand{\HarmCandidates}{19}
\newcommand{\HarmGradings}{36}
\newcommand{\HarmLooSign}{14}
\newcommand{\HarmLooBreaks}{4}
\newcommand{\HelpMM}{0.556}
\newcommand{\HelpText}{0.222}
\newcommand{\HelpDiff}{$+33.3$}
\newcommand{\HelpP}{0.070}
\newcommand{\HelpDiscordant}{8}
\newcommand{\HelpK}{six}
\newcommand{\HelpVerdictK}{one}
\newcommand{\HelpVerdictCandidates}{12}
\newcommand{\HelpFloor}{0.008}
\newcommand{\HelpLooVerdict}{17}
\newcommand{\SigDiff}{$+60.0$}
\newcommand{\SigP}{$4.9 \times 10^{-4}$}
\newcommand{\SigDiscordant}{12}
\newcommand{\SigK}{four}
\newcommand{\SigSignK}{twelve}
\newcommand{\SigHolm}{0.008}
\newcommand{\KillThreshold}{five}
\newcommand{\SmallestK}{one}
\newcommand{\KillState}{did not fire}
\newcommand{\SignificanceNeedsM}{six}
\newcommand{\SweepLow}{0.001}
\newcommand{\SweepHigh}{0.20}
\newcommand{\SweepLevels}{25}
\newcommand{\IncapableSweepMin}{3}
\newcommand{\MedianKSweepMax}{7}
\newcommand{\MedianKSweepMin}{3}
\newcommand{\Incapable}{6}
\newcommand{\AtFloor}{4}
\newcommand{\AtFloorNS}{2}
\newcommand{\AtFloorNSDiscordant}{four}
\newcommand{\SurvivesHolm}{3}
\newcommand{\SignificantRaw}{4}
\newcommand{\Directional}{15}
\newcommand{\FragileMax}{two}
\newcommand{\FragileSign}{4}
\newcommand{\FragileVerdict}{7}
\newcommand{\DecoupledGap}{four}
\newcommand{\DecoupledCount}{9}
\newcommand{\HelpRatio}{six}
\newcommand{\RegradeAnswers}{228}
\newcommand{\RegradeContested}{four}
\newcommand{\RegradeImageFlip}{one}
\newcommand{\RegradeMovingRows}{two}
\newcommand{\RegradeIdenticalRows}{ten}
\newcommand{\ClosedFormMismatches}{0}
\newcommand{\ClosedFormStates}{13,761}
\newcommand{\ClosedFormPairs}{1,201}
\newcommand{\KappaVerInversionsLow}{95}
\newcommand{\KappaVerInversionsHigh}{98}
\newcommand{\KappaVerSigTables}{434}
\newcommand{\SimTables}{840,000}
\newcommand{\SimReps}{10,000}
\newcommand{\FragileGivenSigNEighteen}{0.67}
\newcommand{\FragileGivenSigNThirtyEight}{0.38}
\newcommand{\SpearmanSigNThirtyEight}{$-0.93$}
\newcommand{\NumericHolmLT}{0.066}
\newcommand{\NumericHolmSurvivors}{one}
\newcommand{\LLMKappaQwen}{0.97}
\newcommand{\LLMKappaLlama}{0.91}
\newcommand{\LLMKappaBoth}{0.94}
\newcommand{\RegradedMM}{0.444}
\newcommand{\RegradedDiff}{$+0.0$}
\newcommand{\SummaryHarmStrict}{1}
\newcommand{\SummaryHarmNumeric}{0}
\newcommand{\SummaryHelpStrict}{6}
\newcommand{\SummaryHelpNumeric}{6}
\newcommand{\SummarySigStrict}{4}
\newcommand{\SummarySigNumeric}{2}
\newcommand{\StaticTier}{missing}
\newcommand{\PrimaryTier}{partial}
\newcommand{\SotaTier}{blocked}
\newcommand{\RetrievalN}{38}
\newcommand{\RetrievalCorrect}{38}
\newcommand{\IndexFiles}{0}
\newcommand{\BlockedComponent}{full LLM QA}
\newcommand{\ForgedState}{were not}
\newcommand{\NEvidence}{21}
\newcommand{\EvidenceBytes}{527,252}

%% Draft-only owner placeholder (ruling R15, 2026-09-09). Prints ONLY when the
%% build defines \DRAFTTODO (e.g. pdflatex "\def\DRAFTTODO{1}\input{main}");
%% every packet build leaves it undefined, so no unconfirmed declaration text
%% reaches a PDF. The owner replaces each use with the confirmed sentence.
\providecommand{\ownertodo}[1]{\ifdefined\DRAFTTODO[TODO-OWNER: #1]\fi}

%% Self-contained captions. No cross-references, no section numbers: each
%% caption must be readable on its own if the figure is lifted out.

\newcommand{\CapFlips}{%
How far each conclusion is from being overturned, counted in gradings. Each row
is one paired contrast from a question-answering evaluation in which the same
questions were answered under \ConditionsWord{} retrieval conditions and each
answer was graded right or wrong. Two conclusions are measured per contrast: the direction,
overturned when the leading condition stops leading, and the verdict, overturned
when an exact paired test crosses the conventional threshold of \AlphaLevel{} in
whichever direction it currently sits. Position is the smallest number of
individual gradings that would have to have come out the other way, found by
exhaustive search over reachable tables rather than by a greedy walk, so each is
a minimum and not a bound on one. A cross marks a contrast whose difference is
exactly zero and which therefore has no direction to reverse. The shaded region
is fewer than \KillThreshold{} gradings, the threshold fixed in advance as the
point below which a conclusion counts as resting on individual records. Rows
where the two markers are far apart are contrasts whose direction and whose
verdict are nothing like equally secure. The compact row key is W/L for the
widely/less-documented datasets, I/T/A for image/text/all questions, and M/T/N
for multimodal/text/none retrieval.%
}

\newcommand{\CapRegrade}{%
Of the \RegradeContested{} disagreements, \RegradeImageFlip{} carries the harm contrast. The same
\RegradeAnswers{} answers of a question-answering evaluation were scored by
\Graders{} nested string rules, which agree on all but \RegradeContested{} of
them. The disagreements are reported only in aggregate; question keys, answer
strings, and gold values remain in the private evidence. Panel A gives the
number of answers graded correct under each retrieval condition on the
\SubsetN{} image-dependent questions of the widely-documented-entity set, under
the stored rule (filled marker) and the numeric-tolerant nested rule (open
ring); where the two agree the ring sits on the dot, and only the multimodal
condition moves, by \RegradeImageFlip{} answer. Panel B gives the contrast those
counts produce. Under the stored rule the multimodal condition is behind the
text-only one; under the numeric-tolerant rule the two are exactly level.
The choice between the rules is left open; the point is that the reported
direction is decided by which one is used.%
}

\newcommand{\CapOperating}{%
How often a significant verdict at these sizes rests on one grading. Each point
is the share of simulated paired Bernoulli tables, among those with exact
\(p\le\AlphaLevel{}\), whose verdict distance is one or at most two, plotted
against the number of paired items. The underlying grid is \SimTables{} tables
with \SimReps{} replicates per feasible cell. The thin vertical rules mark the
two design sizes of the evaluation audited here, \SubsetN{} image-dependent
questions and \Questions{} questions. At the \(n\) of this evaluation
the one-grading share is \FragileGivenSigNThirtyEight{}; at the image-subset
size it is \FragileGivenSigNEighteen{}. Both fall as \(n\) grows. The distances
come from the closed form, which exhaustive search confirms on every reachable
table.%
}

\newcommand{\CapFlipSummary}{%
The minimum gradings needed to overturn the three pre-registered claims, shown
under the stored strict grader and the numeric-tolerant nested grader. The first
two rows are direction distances for the widely-documented harm and
less-documented help contrasts; the third is the verdict distance for the
less-documented text-only contrast, and each row carries the same W/L, I/T/A
and M/T/N key as the contrast ladder. The bars of both graders are raw
per-contrast distances; Holm survival is a separate family-level condition and
is reported for the stored grader independently. The numeric-tolerant grader
makes the harm contrast exactly level, so its direction distance changes to
\SummaryHarmNumeric{}; the help direction remains \SummaryHelpNumeric{}, while
the long-tail text verdict changes from \SummarySigStrict{} to
\SummarySigNumeric{}. These are the same fixed paired tables and estimands as
the main analysis, conditioned on grader rather than a new test.%
}

\newcommand{\CapFloor}{%
What each test could have returned, against what it did return. An exact paired
test on binary outcomes is a binomial test on the discordant pairs alone, so the
smallest p-value it can produce is fixed by how many questions the two
conditions disagreed on, before any answer is graded. Each point is one contrast
from a question-answering evaluation; the horizontal position is that floor, the
vertical position is the p-value obtained, and the marker area is the discordant
count. The dashed diagonal is equality, and open markers are contrasts sitting
on it: every discordant pair fell the same way, so the test returned the most
extreme value available to it. A cross marks a contrast on which the two
conditions never disagreed; it also sits on the diagonal, but for a reason that
says nothing about how any answer came out, and it is excluded from the count of
contrasts at their floor. Solid rules mark the conventional threshold of
\AlphaLevel{} on both axes. Points in the shaded region are contrasts whose
floor lies above the threshold: they could not have been called significant
however the answers came out, so their non-significance describes the question
set rather than the conditions compared.%
}

\newcommand{\CapLeaveOneOut}{%
Dropping one question at a time, which is a different question from re-grading
one answer. Each row is one paired contrast from a question-answering
evaluation. Each question in turn is removed and the contrast recomputed on the
remainder; position is the share of those deletions after which the conclusion
still stands, given as a count where it is not unanimous. The direction is
counted as standing if the same condition still leads, and the verdict if an
exact paired test still falls on the same side of the conventional threshold of
\AlphaLevel{}. Most contrasts survive every deletion, which is what makes the
few that do not worth naming: deletion is a blunter instrument than re-grading,
and a conclusion that a single deletion can move was already resting on very
little. Rows use the same W/L, I/T/A and M/T/N key as the flip ladder; a printed
count gives the numerator and denominator whenever the share is not unanimous.%
}

\newcommand{\CapEffectForest}{%
The recorded effect and its exact paired interval, shown for every one of the
\Contrasts{} contrasts. The point is the first-named condition's accuracy minus
the second's in percentage points; the horizontal segment is the exact 95\%
interval conditional on the observed discordant count, obtained by treating the
discordant questions as a binomial split and mapping that split back to the
paired accuracy difference.
The dashed vertical rule is no difference. Each annotation is printed as
``$m/n$'': (m) is the number of discordant questions and (n) is that
contrast's recorded question count. Filled points have a non-zero accuracy
difference; open points are exact zero differences. A zero difference can
still have discordant questions ($m>0$), which is why $m/n$ is printed. W/L identify the widely and
less-documented datasets, I/T/A the image/text/all splits, and M/T/N the
multimodal/text/none retrieval conditions. This is a descriptive re-expression
of the recorded four-cell tables. When $m=0$, no binomial split exists and the
zero-width segment is a conditional marker for the observed zero difference,
not an unconditional 95\% interval. Interval overlap is not used as a
significance rule and no question text or identifier is displayed.%
}

\newcommand{\CapPlane}{%
The plane every grading change moves in. A paired contrast between two
conditions reduces to four counts, and the two an exact test reads are the
discordant pair: the questions the arm alone answered correctly, and the
questions the reference alone did. Regrading one answer moves that pair one unit
along one axis and can do nothing else, so each dot is a table reachable from
the observed one and adjacent dots are one grading apart. Dark dots are tables
at which the conclusion named above the panel no longer holds. The open circle
is the table the evaluation actually produced, and the dashed outline is the set
of tables exactly as many gradings away as the nearest dark dot, cut back to the
tables that exist, so it touches the failure region without entering it. All
three panels are the same lattice over \SubsetN{} questions. The text in the
empty corner of each panel prints the observed discordant counts (b,c) and
whether the named conclusion holds at that table, and the strip above it gives
the minimum distance; both are labels for the observed record, not an
additional state. The panel for the
less-documented direction and the verdict panel that reads the same table are
the same contrast, the same four counts and the same observed point, read as
two different claims: the direction survives \HelpK{} gradings and the
verdict attached to it survives \HelpVerdictK{}. A direction fails on a
half-plane, so its
distance is the lead in discordant pairs; a verdict fails on a region whose
boundary bends, because the test reads how many questions the conditions
disagreed on as well as how the disagreements fell. Two distances measured from
one point to two differently shaped regions have no reason to agree.%
}

\newcommand{\CapDecouple}{%
The two distances a contrast carries, plotted against each other. Each marker is
one paired contrast from a question-answering evaluation, positioned by the
smallest number of individual gradings that would move its verdict across the
conventional threshold of \AlphaLevel{} and by the smallest number that would
reverse its direction. The dashed diagonal is where the two agree and the thin
lines beside it are \DecoupledGap{} gradings either side. Markers in the shaded
lower region are contrasts whose direction is the easier of the two to overturn.
A cross marks a contrast whose difference is exactly zero and which therefore
has no direction to reverse; a small numeral beside a marker is how many
contrasts share that position. Markers appear on both sides of the diagonal,
which is the point: neither claim is systematically the sturdier one, so knowing
one distance does not tell a reader the other, and a results table reporting a
difference and a p-value in one row reports two claims with two different
supports as though they stood together.%
}

\newcommand{\CapAlpha}{%
The same family read against every threshold rather than one. Panel A
counts, for each level on the horizontal axis, how many of the \Contrasts{}
contrasts reach it, how many survive a family-wise correction across all of
them, and how many could not have reached it whatever the answers were, that
last being fixed by how rarely the two conditions disagreed. Panel B
gives the median number of gradings that would move a verdict at that level. The
vertical rule is \AlphaLevel{}, the level the records used and the level every
other figure here is read against. The counts move as the threshold moves and
the shape does not: at every level a substantial share of the contrasts is
outside what the design could resolve, so the observations reported here are
about the evaluation rather than about the conventional choice of threshold. The
sweep recomputes each distance in closed form; the exhaustive search is run at
the recorded level, where the build checks the two against each other.%
}

\newcommand{\CapCellsTable}{%
The four counts each contrast reduces to. Rows are the \Contrasts{} paired
contrasts among \ConditionsWord{} retrieval conditions, over two question sets,
each split into the image-dependent questions, the text-answerable ones, and
both together. For each, the number of questions both conditions answered
correctly, the number the first-named condition alone answered correctly, the
number the second alone did, and the number neither did. The final column is the
sum of the middle two, which is the only part of the table an exact paired test
reads. Every quantity reported anywhere in this paper is a function of these
four numbers and the sample size, so this table is what a reader would need to
recompute any of them.%
}

\newcommand{\CapDesignTable}{%
What a paired comparison can say, as a function of disagreement alone. An exact
paired test on binary outcomes is a binomial test at one half on the questions
the two conditions disagreed about, so both columns in the middle are fixed
before any answer is graded. The second gives the smallest p-value the test can
return on that many discordant pairs, obtained when every one of them falls the
same way. The third gives the smallest majority among them that reaches the
conventional threshold of \AlphaLevel{}, and reads \emph{unreachable} where no
split of that many pairs can. The final column is how many of the \Contrasts{}
contrasts in this evaluation have that discordant count. Nothing in the first
three columns depends on this evaluation: the table is arithmetic on the
binomial and can be read off by anyone comparing two systems on shared items.%
}

\newcommand{\CapGradersTable}{%
The same answers under \Graders{} nested string rules. Every answer of both
question sets was scored against its gold string three ways --- exact match,
substring match, and a rule tolerant of numeric formatting --- and the table
gives the accuracy each rule produces for each retrieval condition, on the
image-dependent and text-answerable splits of both sets. \RegradeIdenticalRows{}
of the twelve rows are identical across the three rules. \RegradeMovingRows{}
move: the multimodal condition on the image-dependent split of the
widely-documented set, and the no-retrieval condition on the text split of the
less-documented set. The table is included so that claim can be checked rather
than accepted: the disagreement between nested rules is visible here as the
only places any number moves.%
}

\newcommand{\CapCandidatesTable}{%
Aggregate counts of single-grading candidates. Where the smallest number of
gradings that would overturn a conclusion is one, the evidence enumerates every
qualifying row, and this table reports how many there are. Question keys,
answer strings, and gold values are withheld from the table; row-level audit
goes through the redacted evidence release that accompanies the manuscript.
Several candidates do not make a conclusion sturdier: they mean several answers
each carry it alone.%
}

\newcommand{\CapContrastsTable}{%
Every paired contrast the evaluation produced, with what it would take to
overturn each. Rows are the \Contrasts{} contrasts among \ConditionsWord{}
retrieval conditions, over two question sets, each split into the image-dependent
questions, the text-answerable ones, and both together. Differences are in
percentage points and are positive when the first condition named leads.
P-values are two-sided exact paired tests; the floor beside each is the smallest
value that test could have returned given how many questions the two conditions
disagreed on, and a p-value equal to its floor means every disagreement fell the
same way. The two flip columns are the smallest number of individual gradings
that would have to have come out otherwise for the direction to reverse and for
the verdict to cross \AlphaLevel{}, each found by exhaustive search. The final
column adjusts the p-values across all \Contrasts{} contrasts as a single
family, the correction the records name as the one that would raise the bar. A
direction column of zero marks a difference of exactly zero, which has no
direction to reverse.%
}

\newcommand{\CapPreregTable}{%
The three conclusions named in advance, and one that was not. Each row is a
statement the evaluation's own records make, followed by the smallest number of
individual gradings that would have to have come out otherwise for it to stop
holding. Where that number is one, the candidate column gives how many distinct
single gradings would each suffice on their own; several candidates do not make
a conclusion sturdier, they mean several answers each carry it alone. The final
column is how many of the single-question deletions leave the conclusion in
place. The first three rows were fixed before the search was written: two are
about which condition leads, the third about whether a test calls a difference
significant. The final row is the verdict on the same contrast as the second
row, separated out because the direction of that difference and the verdict on
it are not the same distance from being overturned.%
}

\newcommand{\CapSecondEvalTable}{%
A labelled broader check, not a replacement for the
\Contrasts{}-contrast evaluation. Each row is one public paired binary
task on a pre-declared model pair (\SecondEvalArmA{} and
\SecondEvalArmB{}), scored from stored item-level correctness rather
than from a new grading. The two flip columns are the same distances
used on the bound evaluation: the smallest number of individual
gradings that would reverse the direction, and the smallest that would
move an exact paired verdict across \AlphaLevel{}. The gap is the
absolute difference of those two counts. Holm is the family adjustment
over the \SecondEvalN{} retained contrasts. Accuracies are omitted on
purpose: this table is not a model ranking. The prestated reading was
that single-record fragility at $n=\SubsetN{}$ does not transport if
every distance is at least \SecondEvalFragileMax{}, and that
direction--verdict disagreement transports if any gap is at least
\DecoupledGap{}.%
}


\title{Count the gradings, not the questions: how far each conclusion of a small
paired evaluation is from being overturned}
\author{Working draft --- author list not yet fixed}
\date{Working draft. Not submitted to any venue.}

\begin{document}
\maketitle

\begin{abstract}
Small paired evaluations report conclusions, and a conclusion is a statement
about how a handful of answers happened to be graded. This paper takes one
completed question-answering evaluation --- \Questions{} questions answered
under \ConditionsWord{} retrieval conditions, \Answers{} graded answers in total --- and
asks of each of its \Contrasts{} recorded contrasts the same arithmetic
question: what is the fewest individual gradings that would have to have come
out the other way for the reported conclusion to stop holding? Because a paired
contrast reduces to four counts and one grading moves a question between two of
them, the answer is found by exhaustive search and is a minimum rather than a
bound. Three conclusions were named in advance. The finding that the multimodal
condition trails the text-only one on image-dependent questions rests on
\HarmK{} grading: \HarmCandidates{} of the \HarmGradings{} individual gradings
on that subset would each on their own erase it. That flip is not hypothetical.
An independent regrade of the same answers by \Graders{} graders finds exactly
\RegradeContested{} grading they disagree about, and applying it as the record
states moves the difference from \HarmDiff{} to \RegradedDiff{} percentage
points. The second conclusion, a \HelpDiff{}-point difference in the other
direction, needs \HelpK{} gradings to reverse --- but its non-significant
verdict needs \HelpVerdictK{}, so the direction of that result is \HelpRatio{}
times better supported than the verdict attached to it. The third, the only
contrast to survive correction for the \Contrasts{} comparisons, needs \SigK{}.
Two further facts about the same rows sharpen the point: \Incapable{} of the
\Contrasts{} contrasts had a smallest attainable p-value above the threshold and
so could not have been called significant however the answers came out, and
\AtFloor{} returned exactly the most extreme value available to them. None of this
adjudicates whether the retrieval conditions differ, and no claim that they do
or do not is made here. What it establishes is narrower and checkable: at this
scale a conclusion is a property of individual gradings, the direction of a
result and the verdict on it are separately fragile and can differ by a factor
of \HelpRatio{}, and reporting either without the count of gradings it rests on
presents an arithmetic accident as a finding.

\end{abstract}

\section{Introduction}

A small evaluation reports that one condition beats another. Underneath that
sentence is a table of questions, each marked right or wrong, and the sentence
is true because of how a few of those marks fell. How few is a question with an
exact answer, and it is almost never asked.

This paper asks it of one completed evaluation. A question set of \Questions{}
items was answered by an instruction model under \ConditionsWord{} retrieval
conditions and graded against gold strings, twice over --- once on widely
documented entities and once on less documented ones --- for \Answers{} graded
answers. The project that produced those runs computed \Contrasts{} paired
contrasts from them and recorded which it considered to have found something.
Nothing here re-runs any of that. The rows are on disk, the recorded contrasts
are on disk, and the only new object is a count: for each conclusion, the
smallest number of individual gradings that would have to have come out the
other way for it to stop holding.

The count is exact rather than estimated. A paired contrast between two
conditions reduces to four numbers --- how many questions both got right, how
many each got right alone, how many neither got right --- and every quantity
computed from it, including the exact test, is a function of those four. One
grading change moves one question from one of the four to an adjacent one. So
the tables reachable by $k$ grading changes are exactly those reachable in $k$
steps, the space is small, and a breadth-first walk returns the true minimum.

Three conclusions were named before the search was written, and the answers are
not alike. The reported finding that the multimodal condition trails the
text-only one on image-dependent questions rests on \HarmK{} grading, and
\HarmCandidates{} of the \HarmGradings{} individual gradings on that subset
would each on their own erase it. The reported \HelpDiff{}-point difference in
the other direction needs \HelpK{} gradings to reverse. The one contrast that
survives correction for multiplicity needs \SigK{}. The prespecified kill for
this study was that all three come out at \KillThreshold{} or more, in which
case the claim that these conclusions turn on individual records would have
failed; it \KillState{}, and the smallest is \SmallestK{}.

Two things make this more than an arithmetic exercise. The first is that the
single grading is not hypothetical. An independent regrade of the same answers
by three graders finds exactly \RegradeContested{} answer they disagree about,
and it is one of the \HarmCandidates{}: applied as the record states it, the
reported difference moves from \HarmDiff{} to \RegradedDiff{} percentage points.
The conclusion and its negation are separated by one adjudication of one string.

The second is that a contrast carries two conclusions, not one, and they are not
equally secure. The \HelpDiff{}-point difference needs \HelpK{} gradings to
change direction, which is comparatively sturdy at this scale. Its verdict ---
that an exact test does not call it significant --- needs \HelpVerdictK{}. On
the same four counts, from the same \SubsetN{} questions, the direction is
\HelpRatio{} times further from being overturned than the significance statement
attached to it. Among the \Directional{} contrasts that have a direction to
reverse at all, the two distances differ by \DecoupledGap{} gradings or more in
\DecoupledCount{} of them. Reporting a contrast as ``not significant''
and reporting it as ``smaller'' are not the same claim and are not equally
supported, and nothing in a conventional results table distinguishes them.

A third observation falls out of the same rows and is arguably the most
uncomfortable. Of the \Contrasts{} contrasts, \Incapable{} have a smallest
attainable p-value above the threshold: the two conditions disagreed on so few
questions that no assignment of those disagreements could have produced a
significant result. Their non-significance is a fact about the question set and
carries no information about the conditions being compared. A further
\AtFloor{} returned exactly the most extreme value available to them, so they
were at the limit of what the design could say.

None of this adjudicates the underlying comparison, and this paper makes no
claim about whether the retrieval conditions differ. The reference benchmark
this line of work would be measured against was never obtained, the tier that
would hold that comparison is recorded as \PrimaryTier{} and the tier above it
as \SotaTier{}, and the counting below does not stand in for either. The
contribution is the count and what it exposes: at this scale a conclusion is a
property of individual gradings, and a results table that prints a difference
and a p-value while omitting how many gradings each rests on has withheld the
part a reader would need in order to know what they are looking at.

\section{Methods}
%% Independent methods file. Not woven into main.tex prose.
%% Every quantity described here resolves to a hash-bound entry in the manifest
%% before it may be cited in Results.

\subsection{What was already on disk}

Nothing in this study was generated for it. No model was run, no answer was
produced, and no question was written. Every quantity below comes from records
written by runs that finished before this analysis began, and each record is
bound to the manuscript by its SHA-256 digest; the analysis code re-hashes each
file before reading it and stops the build on any mismatch.

The evaluation being re-read is small and simple. A question set of
\Questions{} questions was answered by a \Model{} instruction model at
temperature \Temperature{} under \ConditionsWord{} retrieval conditions --- no
retrieval, retrieval over a text-only knowledge graph, and retrieval over the
same graph with image-derived attributes attached --- and each answer was
graded right or wrong against a gold string. The latter two are the
graph-structured form of retrieval-augmented
generation~\cite{lewis2020retrieval,edge2024local} and differ only in whether
the graph carries attributes derived from images. Two such question sets exist, one
built around widely documented entities and one around less documented ones,
giving \Answers{} graded answers in total. Each set is split by a recorded
per-question flag into the \SubsetN{} questions whose answer depends on the
image and the \TextOnlyN{} that can be answered from text alone.

The design is paired: the same question is answered under every condition, so a
contrast between two conditions is decided entirely by the questions the two
conditions disagree on. Everything in this paper follows from that.

\subsection{The four counts a contrast reduces to}

For two conditions over a set of questions, each question falls into exactly one
of four cells: both conditions right, the first right alone, the second right
alone, or neither right. The difference in accuracy, the exact test and its
interval are all functions of those four counts and of nothing else.

Changing one grading --- deciding that one answer under one condition should
have been marked the other way --- moves one question from one cell to an
adjacent one. So the set of tables reachable from the observed one by $k$
grading changes is exactly the set reachable in $k$ steps of that move, and the
question this paper asks has an exact answer rather than an estimate.

\subsection{Two conclusions per contrast, and two distances}

A contrast can carry two different claims, and they are unseated by different
things. The first is a direction: which condition is ahead. The second is a
verdict: whether an exact test calls the difference significant at the
conventional threshold of \AlphaLevel{}. This study measures the distance to
each separately, because its central observation is that they are frequently not
the same distance.

The direction is unseated when the difference reaches zero or reverses. The
verdict is unseated when the p-value crosses the threshold, in whichever
direction it currently sits: a significant result is unseated by becoming
non-significant, and a non-significant one by becoming significant. The second
half of that is not symmetry for its own sake. A verdict of ``no effect'' is
reported and acted on exactly as a verdict of ``effect'' is, so it is entitled
to the same audit.

\subsection{Finding the minimum, rather than bounding it}

The search is exhaustive over reachable states rather than greedy over answers.
Because a contrast is fully described by four non-negative counts summing to
$n$, the reachable tables form a small graph, and a breadth-first walk outward
from the observed table returns the first state satisfying the stopping
condition. That state is at the true minimum distance, so the numbers reported
here are minima and not upper bounds on minima.

Where the minimum is one, the paper enumerates every individual grading that
would suffice on its own, so a reader can see which answers carry a conclusion
rather than being told how many. Several candidates do not make a conclusion
sturdier: it means there are several answers, any one of which carries it alone.

\subsection{What the test could have returned}

An exact test on paired binary data is a binomial test at one half on the
discordant pairs alone~\cite{mcnemar1947note}, which is the standard choice for
comparing two systems on a shared test set~\cite{dror2018hitchhiker}, so the
smallest p-value it can
return is fixed by the number of discordant pairs before any answer is graded.
That floor is reported for every contrast. A contrast whose floor lies above the
threshold could not have been called significant however the answers came out,
and its non-significance is therefore a fact about the question set rather than
about the conditions being compared.

Intervals are Clopper--Pearson on the discordant split mapped onto the paired
difference~\cite{clopper1934use}, matching the construction the records already
carry. Where a contrast has no discordant pairs the test is undefined and its
p-value is reported as one rather than as an error.

\subsection{Deletion, as a separate question}

Beside each flip distance the paper reports a leave-one-out
count~\cite{efron1979bootstrap}: how many of the $n$ single-question deletions
leave the direction, and separately the verdict, as it was. This measures
something different. A flip asks what a different grading of the same question
set would have done; a deletion asks what a different question set would have
done. A conclusion can be robust to one and fragile to the other, and both are
reported for every contrast rather than one standing in for the other.

\subsection{Multiplicity}

The records report \Contrasts{} paired contrasts and state that their p-values
are unadjusted. This paper adds the adjustment the record names as the one that
would raise the bar~\cite{holm1979simple} and reports it beside each raw value.
The adjustment is applied to the whole family of \Contrasts{}, because that is
the family that was computed; no subset of it is treated as though it had been
specified in advance.

\subsection{Checks the build enforces}

Several sentences in this paper are enforced as build-time checks rather than
asserted in prose, and the analysis code stops rather than reporting a number
that would contradict them.

Every one of the \Contrasts{} contrasts is recomputed from the per-answer rows
and checked against the recorded version on three quantities separately: the
number of questions, the difference in accuracy, and the exact p-value. The
discordant split is checked as well as the difference, because a table matching
on the difference alone could produce the same p-value by coincidence. The
recorded verdict is checked to follow from the recorded test rather than being
taken on its word.

The reshaping from per-answer rows into paired tables refuses a question that
was not answered under every condition, refuses a condition that answers the
same question twice, and refuses a question whose subset flag differs between
conditions --- each of which would let a pair be formed across the wrong things.
Each run's own summary is checked to still restate the rows in its own
directory, and the fixed measurement table the project recorded as its result is
checked to still come from those rows.

The regrade record is checked to contain exactly one contested grading, and to
be the one the manuscript names. The value the project recorded as not to be
used is checked to still be recorded that way. The tiers reported here as absent
or incomplete are checked to still be in that state, the blocked component is
checked to still be blocked with no score substituted for it, and the retrieval
result is checked to still carry the index state that keeps it a footnote.

\subsection{What is deliberately not computed}

No answer is regenerated and no grading is decided here; the contested grading
discussed below is one the records already contain, applied as the record states
it, not a judgement made by this paper. The reference benchmark this line of
work would eventually be measured against was never downloaded, the tier that
would hold that comparison is recorded as \PrimaryTier{} and the tier above it
as \SotaTier{}, and no value is imputed for either. The \BlockedComponent{}
component is recorded as blocked on a service that was unreachable, and its
scores \ForgedState{} invented in its place. Nothing in this paper substitutes
for any of them.


\section{How far each conclusion is from being overturned}

Table~\ref{tab:contrasts} gives all \Contrasts{} contrasts and
Figure~\ref{fig:flips} plots the two distances for each. Every recorded
difference and every recorded p-value reproduced from the per-answer rows, and
so did every discordant split; the build would have stopped otherwise, and what
follows is therefore a re-reading of the project's own numbers rather than a
second measurement of them.

The three prespecified conclusions are in Table~\ref{tab:prereg}. Taking them in
turn.

The first is that on the image-dependent questions of the widely-documented set,
the multimodal condition is behind the text-only one: \HarmMM{} against
\HarmText{} over \SubsetN{} questions, a difference of \HarmDiff{} percentage
points on \HarmDiscordant{} discordant pairs, with an exact p-value of
\HarmP{}. The minimum flip set is \HarmK{} grading, and the enumeration is
worth more than the count: \HarmCandidates{} of the \HarmGradings{} individual
gradings on that subset would each, alone, be enough to remove the difference or
reverse it. Slightly more than half of the individual judgements underlying that
subset are each independently sufficient to erase its reported direction.

The second is that on the image-dependent questions of the less-documented set,
the multimodal condition is ahead: \HelpMM{} against \HelpText{}, a difference
of \HelpDiff{} percentage points, exact p-value \HelpP{}, which the records
report as not significant and which this paper does not treat as an established
effect. Its direction needs \HelpK{} gradings to reverse. That is the sturdiest
direction among the three at this subset size, and it is still \HelpK{} answers
out of \SubsetNWord{}.

The third is the only contrast in the study that survives adjustment across the
family of \Contrasts{}: on the text-answerable questions of the less-documented
set, the text-retrieval condition beats no retrieval by \SigDiff{} percentage
points with an exact p-value of \SigP{}, or \SigHolm{} after correction. Its
verdict needs \SigK{} gradings and its direction needs \SigSignK{}. This is what
a comparatively secure result looks like at this scale, and it is also the least
interesting one available: it says that supplying a fact makes a model more
likely to state that fact.

\begin{table}[tbp]
\centering
\caption{\CapPreregTable}
\label{tab:prereg}
\small
% Generated by the figure script from hash-bound evidence. Do not hand-edit.
\begin{tabular}{lrrrr}
\toprule
 & As & Gradings & Single-grading & Deletions \\
Conclusion & measured & to unseat & candidates & survived \\
\midrule
Multimodal is behind text, famous & $-5.6$ & 1 & 19 & 14 of 18 \\
Multimodal is ahead of text, long-tail & $+33.3$ & 6 & --- & 18 of 18 \\
Text beats no retrieval, long-tail & $+60.0$, $4.9 \times 10^{-4}$ & 4 & --- & 20 of 20 \\
\midrule
Same contrast, verdict not direction & 0.070 & 1 & 12 & 17 of 18 \\
\bottomrule
\end{tabular}

\end{table}

\begin{figure}[tbp]
\centering
\includegraphics[width=\linewidth]{fig1_flips.pdf}
\caption{\CapFlips}
\label{fig:flips}
\end{figure}

\subsection{The direction and the verdict are not equally secure}

The fourth row of Table~\ref{tab:prereg} is not one of the three, and it is the
result this paper would keep if it could keep only one. The
\HelpDiff{}-point difference whose direction needs \HelpK{} gradings has a
verdict --- not significant at \AlphaLevel{} --- that needs \HelpVerdictK{}, and
\HelpVerdictCandidates{} distinct single gradings would each be enough to make
it significant. Its floor is \HelpFloor{}, so it was capable of significance and
simply did not reach it.

The same four counts, read two ways, give a claim that \HelpK{} answers could
not shake and a claim that \HelpVerdictK{} answer could. Which of the two a
reader takes away is decided by the conventions of the table it is printed in,
and a table reporting $p = \HelpP{}$ with no further annotation invites the
reading that has the weaker support.

Across the study this is the common case rather than the exception. In
\DecoupledCount{} of the \Directional{} contrasts that have a direction to
reverse, the two distances differ by \DecoupledGap{} gradings or more, and they
differ in both directions: \FragileSign{} contrasts have a direction that
\FragileMax{} gradings or fewer would reverse, \FragileVerdict{} have a verdict
that \FragileMax{} or fewer would flip, and no contrast is in both groups.
There is no general rule that one is sturdier. There is only the
fact that they are different questions, and that reporting one number for both
loses the distinction.

\subsection{The single flip is on record}

Figure~\ref{fig:regrade} makes the first conclusion concrete. The same
\RegradeAnswers{} answers were scored by \Graders{} graders, which agree on all
but \RegradeContested{}. The disagreement is question \RegradeQuestion{} under the
multimodal condition: the model answered \RegradeResponse{} where the gold value
is \RegradeGold{}, which a strict string match marks wrong and a
numeric-tolerant match marks right. It is one of the \HarmCandidates{} gradings
the search identified.

Applying it as the record states moves the multimodal condition from
\HarmMM{} to \RegradedMM{} on those \SubsetN{} questions and the contrast from
\HarmDiff{} to \RegradedDiff{} percentage points. The reported direction does
not weaken; it disappears. Under one grader the multimodal condition is behind,
under the other the two are exactly level, and the two graders differ on how to
read one numeric string.

This paper does not endorse either grader, and the choice between them is a real
methodological question rather than a formality --- whether a tolerance is
appropriate depends on what the question was asking, and that is a judgement
about the item. The point is upstream of the choice. A reported direction that
is decided by an unstated grading convention was never a property of the
conditions being compared, and the way to notice that before publishing is to
count, because the count is what makes a grading convention visible as a
load-bearing part of the result.

That the flip turned up in an independent regrade rather than in the search is
also the reason to trust the search. The two procedures had nothing to do with
each other: one enumerated gradings that would change a conclusion, the other
re-scored answers under different string-matching rules, and they met at the
same answer.

\begin{figure}[tbp]
\centering
\includegraphics[width=\linewidth]{fig2_regrade.pdf}
\caption{\CapRegrade}
\label{fig:regrade}
\end{figure}

\section{What the tests could have returned}

Figure~\ref{fig:floor} reports a quantity that is available before any answer is
graded. An exact paired test is a binomial test on the discordant pairs
alone~\cite{mcnemar1947note}, so the smallest p-value it can produce is fixed by
how many questions the two conditions disagreed on. With three discordant pairs
the most extreme result attainable is a two-sided p-value of one quarter, and no
grading of those three answers can do better.

Of the \Contrasts{} contrasts, \Incapable{} have a floor above \AlphaLevel{}.
They were reported as non-significant, and they could not have been reported as
anything else. For those contrasts the phrase carries no information about the
conditions being compared: it is a restatement of how rarely the two conditions
disagreed, which is a property of the question set.

A second reading of the same figure is the \AtFloor{} contrasts sitting on the
diagonal, where every discordant pair fell the same way and the test returned
the most extreme value available to it. Of those, \AtFloorNS{} are
non-significant regardless. A contrast in which one condition won every single
question the two conditions disagreed about, and which is nonetheless reported
as no evidence of a difference, is the clearest available demonstration that a
threshold verdict at this scale is a statement about $n$ before it is a
statement about the conditions.

Of the \Contrasts{} contrasts, \SignificantRaw{} reach the threshold unadjusted
and \SurvivesHolm{} survive correction across the family~\cite{holm1979simple}.
The records state that their p-values are unadjusted and that a correction would
raise the bar; this applies it. The survivors are all differences against no
retrieval at all, and none of them is a comparison between the two retrieval
conditions, which is the comparison the evaluation was built to make.

\begin{figure}[tbp]
\centering
\includegraphics[width=0.62\linewidth]{fig3_floor.pdf}
\caption{\CapFloor}
\label{fig:floor}
\end{figure}

\begin{table}[tbp]
\centering
\caption{\CapContrastsTable}
\label{tab:contrasts}
\scriptsize
\setlength{\tabcolsep}{4.5pt}
% Generated by the figure script from hash-bound evidence. Do not hand-edit.
\begin{tabular}{llrrrrrrrr}
\toprule
& & & & & & & \multicolumn{2}{c}{Gradings to unseat} & \\
\cmidrule(lr){8-9}
Questions & Contrast & $n$ & Disc. & Diff. (pt) & $p$ & Floor & Dir. & Verd. & Holm \\
\midrule
Famous, image & multimodal vs text & 18 & 7 & $-5.6$ & 1.000 & 0.016 & 1 & 5 & 1.000 \\
Famous, image & text vs none & 18 & 2 & $+0.0$ & 1.000 & 0.500 & 0 & 6 & 1.000 \\
Famous, image & multimodal vs none & 18 & 7 & $-5.6$ & 1.000 & 0.016 & 1 & 5 & 1.000 \\
Famous, text & multimodal vs text & 20 & 0 & $+0.0$ & 1.000 & 1.000 & 0 & 6 & 1.000 \\
Famous, text & text vs none & 20 & 4 & $+20.0$ & 0.125 & 0.125 & 4 & 2 & 1.000 \\
Famous, text & multimodal vs none & 20 & 4 & $+20.0$ & 0.125 & 0.125 & 4 & 2 & 1.000 \\
Famous, all & multimodal vs text & 38 & 7 & $-2.6$ & 1.000 & 0.016 & 1 & 5 & 1.000 \\
Famous, all & text vs none & 38 & 6 & $+10.5$ & 0.219 & 0.031 & 4 & 2 & 1.000 \\
Famous, all & multimodal vs none & 38 & 11 & $+7.9$ & 0.549 & $9.8 \times 10^{-4}$ & 3 & 4 & 1.000 \\
Long-tail, image & multimodal vs text & 18 & 8 & $+33.3$ & 0.070 & 0.008 & 6 & 1 & 0.984 \\
Long-tail, image & text vs none & 18 & 4 & $-11.1$ & 0.625 & 0.125 & 2 & 4 & 1.000 \\
Long-tail, image & multimodal vs none & 18 & 8 & $+22.2$ & 0.289 & 0.008 & 4 & 2 & 1.000 \\
Long-tail, text & multimodal vs text & 20 & 0 & $+0.0$ & 1.000 & 1.000 & 0 & 6 & 1.000 \\
Long-tail, text & text vs none & 20 & 12 & $+60.0$ & $4.9 \times 10^{-4}$ & $4.9 \times 10^{-4}$ & 12 & 4 & 0.008 \\
Long-tail, text & multimodal vs none & 20 & 12 & $+60.0$ & $4.9 \times 10^{-4}$ & $4.9 \times 10^{-4}$ & 12 & 4 & 0.008 \\
Long-tail, all & multimodal vs text & 38 & 8 & $+15.8$ & 0.070 & 0.008 & 6 & 1 & 0.984 \\
Long-tail, all & text vs none & 38 & 16 & $+26.3$ & 0.021 & $3.1 \times 10^{-5}$ & 10 & 2 & 0.319 \\
Long-tail, all & multimodal vs none & 38 & 20 & $+42.1$ & $4.0 \times 10^{-4}$ & $1.9 \times 10^{-6}$ & 16 & 6 & 0.007 \\
\bottomrule
\end{tabular}

\end{table}

\section{Deleting a question is not re-grading an answer}

Figure~\ref{fig:loo} reports the leave-one-out count beside the flip distance,
and the two disagree in an instructive way. Most contrasts survive every single
deletion: the direction and the verdict are unchanged in all $n$ recomputations.
Deletion is a blunt instrument here, because removing a concordant question
changes neither the discordant split nor the sign, and most questions are
concordant.

The exceptions are the contrasts already identified. The direction of the first
conclusion survives \HarmLooSign{} of \SubsetN{} deletions, so \HarmLooBreaks{}
individual questions each carry it by their presence as well as by their
grading. The verdict on the second survives \HelpLooVerdict{} of \SubsetN{}.

The two measurements answer different counterfactuals and should be read as
such. A flip asks what a different adjudication of the same evaluation would
have concluded, which is a question about grading policy and about annotator
disagreement~\cite{northcutt2021pervasive}. A deletion asks what a slightly
different question set would have concluded, which is a question about sampling.
An evaluation can be exposed on one and not the other, and reporting only the
one that happens to look better is a degree of freedom worth closing in
advance~\cite{simmons2011false}.

\begin{figure}[tbp]
\centering
\includegraphics[width=\linewidth]{fig4_loo.pdf}
\caption{\CapLeaveOneOut}
\label{fig:loo}
\end{figure}

\section{Discussion}

\paragraph{The count is old; the target is not.} Counting how many outcomes
would have to change to overturn a result is the fragility
index~\cite{feinstein1990unit,walsh2014statistical}, introduced for randomised
trials, where it revealed that a large share of results in the clinical
literature turn on a handful of events --- one mechanism behind the more general
observation that small studies produce unreliable
findings~\cite{ioannidis2005why}. The construction transfers to machine
learning evaluation almost unchanged, and the reasons to want it are stronger
here: an evaluation item is graded by a rule that someone wrote, benchmark labels
are known to carry errors at rates that change published
rankings~\cite{northcutt2021pervasive}, and the evaluation sets that appear in
practice are frequently small enough that the count is in single
digits~\cite{card2020little}. The main criticism of the index in its original
setting is that it is a monotone function of the p-value and so adds
nothing~\cite{carter2017fragility}. That criticism has less force here, for a
reason this study demonstrates rather than argues: the count is being applied to
the direction of a difference as well as to the verdict on it, the two are not
monotone in each other, and \DecoupledCount{} of the \Directional{} contrasts
that have a direction to reverse separate them by \DecoupledGap{} gradings or
more.

\paragraph{Two conclusions per contrast, and only one of them gets audited.}
Significance is routinely questioned and directions are routinely believed. This
study found the opposite pattern at least as often. The clearest case is the
contrast whose direction \HelpK{} gradings could not reverse and whose verdict
\HelpVerdictK{} grading would flip, but the general point is that these are separate claims with
separate supports, and a results table reports them in a single row as though
they stood or fell together. The recommendation implied is small: print the two
counts. They cost nothing to compute, they are exact, and they say what a
confidence interval says while remaining in the units the reader can inspect ---
answers.

\paragraph{Non-significance is not a finding when significance was
unreachable.} Of the \Contrasts{} contrasts here, \Incapable{} could not have
been significant at any grading of their answers. Reporting them as null results
is not conservative; it attributes to the conditions a property of the question
set. The floor is one line of arithmetic from the discordant count and belongs
beside any negative result from a paired design, in the same way that a power
calculation belongs beside a negative result from an unpaired
one~\cite{card2020little,amrhein2019scientists}. The stronger version of the
point is the contrast in which every disagreement fell one way and the verdict
was still negative: unanimity among the informative questions was not enough,
because there were \AtFloorNSDiscordant{} of them.

\paragraph{A grading convention is a researcher degree of freedom.} String
matching against gold answers looks like a mechanical step and is not one.
Strict matching, substring matching and numeric tolerance are three defensible
rules, they disagree about one answer here, and that one answer decides a
reported direction. The literature on undisclosed analytic
flexibility~\cite{simmons2011false,gelman2014statistical} is usually read as
being about model selection and outlier handling; grading rules belong on the
same list, and in small evaluations they may dominate it. The regrade in this
study is the practice that made the point visible, and it cost one pass over the
stored answers with no model in the loop.

\paragraph{Report the count with the claim.} The concrete proposal is one column.
Beside each reported contrast, print the smallest number of gradings that would
overturn its direction and the smallest that would move its verdict. Both are
exhaustively computable at these sizes. A reader seeing \HarmK{} in that column
knows immediately what kind of object they are being shown, without needing to
reconstruct the discordant table from an accuracy and a sample size --- which is
what it currently takes, and which is why the question is rarely asked. This is
the same instinct behind reporting effect sizes and intervals rather than
verdicts~\cite{wasserstein2016asa,amrhein2019scientists}, expressed in the unit
that a small evaluation is actually made of.

\paragraph{What this does not fix.} Counting gradings makes fragility visible
and does nothing about its cause. An evaluation whose conclusions rest on single
answers needs more questions, not better reporting of how few it
had~\cite{card2020little,bowman2021will}. The count also says nothing about
whether the question set measures what it claims to measure, which at these
sizes is the larger risk~\cite{raji2021everything}. Its value is narrower: it is
cheap, it is exact, it applies to records that already exist, and it prevents
one specific error, which is publishing an arithmetic accident in the register
of a finding.

\section{What this study does not claim}

No claim is made that adding image-derived attributes to retrieval helps or
harms. The contrast on the less-documented set is \HelpDiff{} percentage points
with an exact p-value of \HelpP{}, which does not reach the threshold and is not
treated here as an established effect; the contrast on the widely-documented set
is \HarmDiff{} points and is overturned by \HarmK{} grading. Both are reported
as objects whose fragility is being measured, and neither is offered as
evidence about the underlying conditions.

No claim is made about performance on any reference benchmark. The benchmark
this line of work would be measured against was never obtained, the tier that
would hold that comparison is recorded as \PrimaryTier{} and the tier above it
as \SotaTier{}, and no value is imputed for either. No result from any other
system is introduced for comparison.

No retrieval result here is a benchmark accuracy. The retrieval step resolves
\RetrievalCorrect{} of \RetrievalN{} questions to the right entity on each
question set, and that figure is recorded as a footnote on an easy
small-graph lookup with \IndexFiles{} files in the index it names. It is
reported here only to establish that the conditions differ in what they
retrieved rather than in whether retrieval worked at all.

No score is supplied for the blocked component. The \BlockedComponent{}
component is recorded as blocked on a service that was unreachable, its scores
\ForgedState{} invented in its place, and nothing here substitutes for them. The
tier that would hold the full comparison remains \PrimaryTier{} and this paper
does not present its own counting as progress toward completing it.

No use is made of the value the project recorded as not to be used. An earlier
run of the same subset under a different caption source produced a much higher
figure for the multimodal condition; the records mark it as unused and name the
real-caption values as the result, and this paper quotes the real-caption values
only. The excluded figure appears nowhere in this manuscript.

No grading is adjudicated here. The contested answer is contested in the records
already, both readings of it are the records' own, and this paper reports what
each implies rather than deciding between them. Nothing was re-graded, re-run or
re-scored for this study.

No claim is made that a minimum flip set of one means a conclusion is false. It
means the conclusion is supported by one grading. A true effect measured on
eighteen questions can also rest on one grading, and the count does not
distinguish those cases; it only establishes what would be needed to distinguish
them, which is more questions.

\section{Conclusion}

One completed evaluation of \Questions{} questions under \ConditionsWord{} retrieval
conditions was re-read by counting rather than by re-running. For each of its
\Contrasts{} paired contrasts, the smallest number of individual gradings that
would overturn the direction, and separately the verdict, was found by
exhaustive search over the reachable tables.

The three conclusions named in advance came out at \HarmK{}, \HelpK{} and
\SigK{} gradings, so the prespecified kill \KillState{}. The smallest,
\SmallestK{}, belongs to the reported finding that the multimodal condition
trails the text-only one: \HarmCandidates{} of the \HarmGradings{} gradings on
that subset would each erase it alone, and an independent regrade of the same
answers disagrees about exactly one of them, which moves the difference from
\HarmDiff{} to \RegradedDiff{} percentage points. The direction of the opposing
\HelpDiff{}-point difference needs \HelpK{} gradings while the verdict attached
to it needs \HelpVerdictK{}, and across the study those two distances differ by
\DecoupledGap{} or more in \DecoupledCount{} of the \Directional{} contrasts
that have a direction at all. No grading of the answers could have made
\Incapable{} of the contrasts significant.

The recommendation is one column in a results table: beside each contrast, the
number of gradings that would overturn its direction and the number that would
move its verdict. It is exactly computable, it costs nothing at these sizes, and
it converts a question a reader cannot currently ask into a number they can
read. When the count is large, nothing has been lost by printing it. When it is
\HarmK{}, what has been learned is that the row is reporting an adjudication of
one answer, and that is worth knowing before the table is published rather than
after.

\section*{Data and code availability}

The per-answer records for both question sets, the summaries each run wrote for
itself, the recorded paired contrasts with their exact tests and intervals, the
independent regrade under three graders, the fixed measurement table, the
recorded state of each comparison tier including those that are absent or
blocked, and the analysis code that produces every number in this manuscript
accompany the manuscript source. Each artifact is recorded with its SHA-256
digest in a manifest, and the analysis code verifies those digests against the
bytes on disk before reading them, so the build fails rather than silently
reporting stale numbers. The knowledge graph and question sets are small
local artifacts distributed with the records; no third-party benchmark is
redistributed, and the model used to produce the answers is a third-party
release identified by name rather than copied.
%% Generated by scripts/release_targets.py from the release ledger.
%% Do not edit: an identifier typed here would not be checked against
%% anything, which is exactly the failure the ledger exists to prevent.
%% Ledger state: github=CREATED, zenodo=DEPOSITED
That material is deposited under the persistent identifier \href{https://doi.org/10.5281/zenodo.22647028}{10.5281/zenodo.22647028}, and the development repository is at \url{https://github.com/PeterPonyu/minimum-flip-set-audit}.


\section*{Ethics statement}

This study is a secondary analysis of numeric records produced by earlier
computational runs. No human subjects are involved, no personal data are
processed, and no new model output was generated. The questions concern publicly
documented landmarks and comparable entities. One consideration bears directly
on the subject matter: the paper's central observation is that a reported
difference between two systems can be decided by a single grading decision, and
the same is true of evaluations used to support claims about deployed systems,
where a fragile result may be read as a safety or capability assurance.
Reporting the count is offered partly for that reason.

\section*{Competing interests}

Not yet declared. This is a working draft and the declaration is outstanding.

\section*{Funding and author contributions}

Not yet declared. This is a working draft; the author list is not fixed and
contributions are therefore not yet attributable.

\section*{Declaration of generative AI use}

The analysis code, the figure code and the drafting of this manuscript were
produced with substantial assistance from large language model agents working
under human direction. Two safeguards bear directly on what is reported. Every
number printed above is emitted programmatically from artifacts bound by
cryptographic digest rather than transcribed by a model into prose. And several
of the paper's structural claims --- that every recorded difference, discordant
split and exact p-value reproduces from the per-answer rows, that each run's
summary still restates the rows in its own directory, that the regrade record
contains exactly one contested grading and that it is the one named here, and
that the tiers reported as absent or blocked are still in that state with no
score substituted --- are enforced as build-time checks that stop the build
rather than as sentences. The answers being analysed were themselves produced by
a language model in the earlier runs this study re-reads; none was regenerated
here. Responsibility for the content rests with the authors.

\bibliographystyle{plain}
\bibliography{references}

\end{document}
");
%% every packet build leaves it undefined, so no unconfirmed declaration text
%% reaches a PDF. The owner replaces each use with the confirmed sentence.
\providecommand{\ownertodo}[1]{\ifdefined\DRAFTTODO[TODO-OWNER: #1]\fi}

%% Self-contained captions. No cross-references, no section numbers: each
%% caption must be readable on its own if the figure is lifted out.

\newcommand{\CapFlips}{%
How far each conclusion is from being overturned, counted in gradings. Each row
is one paired contrast from a question-answering evaluation in which the same
questions were answered under \ConditionsWord{} retrieval conditions and each
answer was graded right or wrong. Two conclusions are measured per contrast: the direction,
overturned when the leading condition stops leading, and the verdict, overturned
when an exact paired test crosses the conventional threshold of \AlphaLevel{} in
whichever direction it currently sits. Position is the smallest number of
individual gradings that would have to have come out the other way, found by
exhaustive search over reachable tables rather than by a greedy walk, so each is
a minimum and not a bound on one. A cross marks a contrast whose difference is
exactly zero and which therefore has no direction to reverse. The shaded region
is fewer than \KillThreshold{} gradings, the threshold fixed in advance as the
point below which a conclusion counts as resting on individual records. Rows
where the two markers are far apart are contrasts whose direction and whose
verdict are nothing like equally secure. The compact row key is W/L for the
widely/less-documented datasets, I/T/A for image/text/all questions, and M/T/N
for multimodal/text/none retrieval.%
}

\newcommand{\CapRegrade}{%
Of the \RegradeContested{} disagreements, \RegradeImageFlip{} carries the harm contrast. The same
\RegradeAnswers{} answers of a question-answering evaluation were scored by
\Graders{} nested string rules, which agree on all but \RegradeContested{} of
them. The disagreements are reported only in aggregate; question keys, answer
strings, and gold values remain in the private evidence. Panel A gives the
number of answers graded correct under each retrieval condition on the
\SubsetN{} image-dependent questions of the widely-documented-entity set, under
the stored rule (filled marker) and the numeric-tolerant nested rule (open
ring); where the two agree the ring sits on the dot, and only the multimodal
condition moves, by \RegradeImageFlip{} answer. Panel B gives the contrast those
counts produce. Under the stored rule the multimodal condition is behind the
text-only one; under the numeric-tolerant rule the two are exactly level.
The choice between the rules is left open; the point is that the reported
direction is decided by which one is used.%
}

\newcommand{\CapOperating}{%
How often a significant verdict at these sizes rests on one grading. Each point
is the share of simulated paired Bernoulli tables, among those with exact
\(p\le\AlphaLevel{}\), whose verdict distance is one or at most two, plotted
against the number of paired items. The underlying grid is \SimTables{} tables
with \SimReps{} replicates per feasible cell. The thin vertical rules mark the
two design sizes of the evaluation audited here, \SubsetN{} image-dependent
questions and \Questions{} questions. At the \(n\) of this evaluation
the one-grading share is \FragileGivenSigNThirtyEight{}; at the image-subset
size it is \FragileGivenSigNEighteen{}. Both fall as \(n\) grows. The distances
come from the closed form, which exhaustive search confirms on every reachable
table.%
}

\newcommand{\CapFlipSummary}{%
The minimum gradings needed to overturn the three pre-registered claims, shown
under the stored strict grader and the numeric-tolerant nested grader. The first
two rows are direction distances for the widely-documented harm and
less-documented help contrasts; the third is the verdict distance for the
less-documented text-only contrast, and each row carries the same W/L, I/T/A
and M/T/N key as the contrast ladder. The bars of both graders are raw
per-contrast distances; Holm survival is a separate family-level condition and
is reported for the stored grader independently. The numeric-tolerant grader
makes the harm contrast exactly level, so its direction distance changes to
\SummaryHarmNumeric{}; the help direction remains \SummaryHelpNumeric{}, while
the long-tail text verdict changes from \SummarySigStrict{} to
\SummarySigNumeric{}. These are the same fixed paired tables and estimands as
the main analysis, conditioned on grader rather than a new test.%
}

\newcommand{\CapFloor}{%
What each test could have returned, against what it did return. An exact paired
test on binary outcomes is a binomial test on the discordant pairs alone, so the
smallest p-value it can produce is fixed by how many questions the two
conditions disagreed on, before any answer is graded. Each point is one contrast
from a question-answering evaluation; the horizontal position is that floor, the
vertical position is the p-value obtained, and the marker area is the discordant
count. The dashed diagonal is equality, and open markers are contrasts sitting
on it: every discordant pair fell the same way, so the test returned the most
extreme value available to it. A cross marks a contrast on which the two
conditions never disagreed; it also sits on the diagonal, but for a reason that
says nothing about how any answer came out, and it is excluded from the count of
contrasts at their floor. Solid rules mark the conventional threshold of
\AlphaLevel{} on both axes. Points in the shaded region are contrasts whose
floor lies above the threshold: they could not have been called significant
however the answers came out, so their non-significance describes the question
set rather than the conditions compared.%
}

\newcommand{\CapLeaveOneOut}{%
Dropping one question at a time, which is a different question from re-grading
one answer. Each row is one paired contrast from a question-answering
evaluation. Each question in turn is removed and the contrast recomputed on the
remainder; position is the share of those deletions after which the conclusion
still stands, given as a count where it is not unanimous. The direction is
counted as standing if the same condition still leads, and the verdict if an
exact paired test still falls on the same side of the conventional threshold of
\AlphaLevel{}. Most contrasts survive every deletion, which is what makes the
few that do not worth naming: deletion is a blunter instrument than re-grading,
and a conclusion that a single deletion can move was already resting on very
little. Rows use the same W/L, I/T/A and M/T/N key as the flip ladder; a printed
count gives the numerator and denominator whenever the share is not unanimous.%
}

\newcommand{\CapEffectForest}{%
The recorded effect and its exact paired interval, shown for every one of the
\Contrasts{} contrasts. The point is the first-named condition's accuracy minus
the second's in percentage points; the horizontal segment is the exact 95\%
interval conditional on the observed discordant count, obtained by treating the
discordant questions as a binomial split and mapping that split back to the
paired accuracy difference.
The dashed vertical rule is no difference. Each annotation is printed as
``$m/n$'': (m) is the number of discordant questions and (n) is that
contrast's recorded question count. Filled points have a non-zero accuracy
difference; open points are exact zero differences. A zero difference can
still have discordant questions ($m>0$), which is why $m/n$ is printed. W/L identify the widely and
less-documented datasets, I/T/A the image/text/all splits, and M/T/N the
multimodal/text/none retrieval conditions. This is a descriptive re-expression
of the recorded four-cell tables. When $m=0$, no binomial split exists and the
zero-width segment is a conditional marker for the observed zero difference,
not an unconditional 95\% interval. Interval overlap is not used as a
significance rule and no question text or identifier is displayed.%
}

\newcommand{\CapPlane}{%
The plane every grading change moves in. A paired contrast between two
conditions reduces to four counts, and the two an exact test reads are the
discordant pair: the questions the arm alone answered correctly, and the
questions the reference alone did. Regrading one answer moves that pair one unit
along one axis and can do nothing else, so each dot is a table reachable from
the observed one and adjacent dots are one grading apart. Dark dots are tables
at which the conclusion named above the panel no longer holds. The open circle
is the table the evaluation actually produced, and the dashed outline is the set
of tables exactly as many gradings away as the nearest dark dot, cut back to the
tables that exist, so it touches the failure region without entering it. All
three panels are the same lattice over \SubsetN{} questions. The text in the
empty corner of each panel prints the observed discordant counts (b,c) and
whether the named conclusion holds at that table, and the strip above it gives
the minimum distance; both are labels for the observed record, not an
additional state. The panel for the
less-documented direction and the verdict panel that reads the same table are
the same contrast, the same four counts and the same observed point, read as
two different claims: the direction survives \HelpK{} gradings and the
verdict attached to it survives \HelpVerdictK{}. A direction fails on a
half-plane, so its
distance is the lead in discordant pairs; a verdict fails on a region whose
boundary bends, because the test reads how many questions the conditions
disagreed on as well as how the disagreements fell. Two distances measured from
one point to two differently shaped regions have no reason to agree.%
}

\newcommand{\CapDecouple}{%
The two distances a contrast carries, plotted against each other. Each marker is
one paired contrast from a question-answering evaluation, positioned by the
smallest number of individual gradings that would move its verdict across the
conventional threshold of \AlphaLevel{} and by the smallest number that would
reverse its direction. The dashed diagonal is where the two agree and the thin
lines beside it are \DecoupledGap{} gradings either side. Markers in the shaded
lower region are contrasts whose direction is the easier of the two to overturn.
A cross marks a contrast whose difference is exactly zero and which therefore
has no direction to reverse; a small numeral beside a marker is how many
contrasts share that position. Markers appear on both sides of the diagonal,
which is the point: neither claim is systematically the sturdier one, so knowing
one distance does not tell a reader the other, and a results table reporting a
difference and a p-value in one row reports two claims with two different
supports as though they stood together.%
}

\newcommand{\CapAlpha}{%
The same family read against every threshold rather than one. Panel A
counts, for each level on the horizontal axis, how many of the \Contrasts{}
contrasts reach it, how many survive a family-wise correction across all of
them, and how many could not have reached it whatever the answers were, that
last being fixed by how rarely the two conditions disagreed. Panel B
gives the median number of gradings that would move a verdict at that level. The
vertical rule is \AlphaLevel{}, the level the records used and the level every
other figure here is read against. The counts move as the threshold moves and
the shape does not: at every level a substantial share of the contrasts is
outside what the design could resolve, so the observations reported here are
about the evaluation rather than about the conventional choice of threshold. The
sweep recomputes each distance in closed form; the exhaustive search is run at
the recorded level, where the build checks the two against each other.%
}

\newcommand{\CapCellsTable}{%
The four counts each contrast reduces to. Rows are the \Contrasts{} paired
contrasts among \ConditionsWord{} retrieval conditions, over two question sets,
each split into the image-dependent questions, the text-answerable ones, and
both together. For each, the number of questions both conditions answered
correctly, the number the first-named condition alone answered correctly, the
number the second alone did, and the number neither did. The final column is the
sum of the middle two, which is the only part of the table an exact paired test
reads. Every quantity reported anywhere in this paper is a function of these
four numbers and the sample size, so this table is what a reader would need to
recompute any of them.%
}

\newcommand{\CapDesignTable}{%
What a paired comparison can say, as a function of disagreement alone. An exact
paired test on binary outcomes is a binomial test at one half on the questions
the two conditions disagreed about, so both columns in the middle are fixed
before any answer is graded. The second gives the smallest p-value the test can
return on that many discordant pairs, obtained when every one of them falls the
same way. The third gives the smallest majority among them that reaches the
conventional threshold of \AlphaLevel{}, and reads \emph{unreachable} where no
split of that many pairs can. The final column is how many of the \Contrasts{}
contrasts in this evaluation have that discordant count. Nothing in the first
three columns depends on this evaluation: the table is arithmetic on the
binomial and can be read off by anyone comparing two systems on shared items.%
}

\newcommand{\CapGradersTable}{%
The same answers under \Graders{} nested string rules. Every answer of both
question sets was scored against its gold string three ways --- exact match,
substring match, and a rule tolerant of numeric formatting --- and the table
gives the accuracy each rule produces for each retrieval condition, on the
image-dependent and text-answerable splits of both sets. \RegradeIdenticalRows{}
of the twelve rows are identical across the three rules. \RegradeMovingRows{}
move: the multimodal condition on the image-dependent split of the
widely-documented set, and the no-retrieval condition on the text split of the
less-documented set. The table is included so that claim can be checked rather
than accepted: the disagreement between nested rules is visible here as the
only places any number moves.%
}

\newcommand{\CapCandidatesTable}{%
Aggregate counts of single-grading candidates. Where the smallest number of
gradings that would overturn a conclusion is one, the evidence enumerates every
qualifying row, and this table reports how many there are. Question keys,
answer strings, and gold values are withheld from the table; row-level audit
goes through the redacted evidence release that accompanies the manuscript.
Several candidates do not make a conclusion sturdier: they mean several answers
each carry it alone.%
}

\newcommand{\CapContrastsTable}{%
Every paired contrast the evaluation produced, with what it would take to
overturn each. Rows are the \Contrasts{} contrasts among \ConditionsWord{}
retrieval conditions, over two question sets, each split into the image-dependent
questions, the text-answerable ones, and both together. Differences are in
percentage points and are positive when the first condition named leads.
P-values are two-sided exact paired tests; the floor beside each is the smallest
value that test could have returned given how many questions the two conditions
disagreed on, and a p-value equal to its floor means every disagreement fell the
same way. The two flip columns are the smallest number of individual gradings
that would have to have come out otherwise for the direction to reverse and for
the verdict to cross \AlphaLevel{}, each found by exhaustive search. The final
column adjusts the p-values across all \Contrasts{} contrasts as a single
family, the correction the records name as the one that would raise the bar. A
direction column of zero marks a difference of exactly zero, which has no
direction to reverse.%
}

\newcommand{\CapPreregTable}{%
The three conclusions named in advance, and one that was not. Each row is a
statement the evaluation's own records make, followed by the smallest number of
individual gradings that would have to have come out otherwise for it to stop
holding. Where that number is one, the candidate column gives how many distinct
single gradings would each suffice on their own; several candidates do not make
a conclusion sturdier, they mean several answers each carry it alone. The final
column is how many of the single-question deletions leave the conclusion in
place. The first three rows were fixed before the search was written: two are
about which condition leads, the third about whether a test calls a difference
significant. The final row is the verdict on the same contrast as the second
row, separated out because the direction of that difference and the verdict on
it are not the same distance from being overturned.%
}

\newcommand{\CapSecondEvalTable}{%
A labelled broader check, not a replacement for the
\Contrasts{}-contrast evaluation. Each row is one public paired binary
task on a pre-declared model pair (\SecondEvalArmA{} and
\SecondEvalArmB{}), scored from stored item-level correctness rather
than from a new grading. The two flip columns are the same distances
used on the bound evaluation: the smallest number of individual
gradings that would reverse the direction, and the smallest that would
move an exact paired verdict across \AlphaLevel{}. The gap is the
absolute difference of those two counts. Holm is the family adjustment
over the \SecondEvalN{} retained contrasts. Accuracies are omitted on
purpose: this table is not a model ranking. The prestated reading was
that single-record fragility at $n=\SubsetN{}$ does not transport if
every distance is at least \SecondEvalFragileMax{}, and that
direction--verdict disagreement transports if any gap is at least
\DecoupledGap{}.%
}


\title{Count the gradings, not the questions: how far each conclusion of a small
paired evaluation is from being overturned}
\author{Working draft --- author list not yet fixed}
\date{Working draft. Not submitted to any venue.}

\begin{document}
\maketitle

\begin{abstract}
Small paired evaluations report conclusions, and a conclusion is a statement
about how a handful of answers happened to be graded. This paper takes one
completed question-answering evaluation --- \Questions{} questions answered
under \ConditionsWord{} retrieval conditions, \Answers{} graded answers in total --- and
asks of each of its \Contrasts{} recorded contrasts the same arithmetic
question: what is the fewest individual gradings that would have to have come
out the other way for the reported conclusion to stop holding? Because a paired
contrast reduces to four counts and one grading moves a question between two of
them, the answer is found by exhaustive search and is a minimum rather than a
bound. Three conclusions were named in advance. The finding that the multimodal
condition trails the text-only one on image-dependent questions rests on
\HarmK{} grading: \HarmCandidates{} of the \HarmGradings{} individual gradings
on that subset would each on their own erase it. That flip is not hypothetical.
An independent regrade of the same answers by \Graders{} graders finds exactly
\RegradeContested{} grading they disagree about, and applying it as the record
states moves the difference from \HarmDiff{} to \RegradedDiff{} percentage
points. The second conclusion, a \HelpDiff{}-point difference in the other
direction, needs \HelpK{} gradings to reverse --- but its non-significant
verdict needs \HelpVerdictK{}, so the direction of that result is \HelpRatio{}
times better supported than the verdict attached to it. The third, the only
contrast to survive correction for the \Contrasts{} comparisons, needs \SigK{}.
Two further facts about the same rows sharpen the point: \Incapable{} of the
\Contrasts{} contrasts had a smallest attainable p-value above the threshold and
so could not have been called significant however the answers came out, and
\AtFloor{} returned exactly the most extreme value available to them. None of this
adjudicates whether the retrieval conditions differ, and no claim that they do
or do not is made here. What it establishes is narrower and checkable: at this
scale a conclusion is a property of individual gradings, the direction of a
result and the verdict on it are separately fragile and can differ by a factor
of \HelpRatio{}, and reporting either without the count of gradings it rests on
presents an arithmetic accident as a finding.

\end{abstract}

\section{Introduction}

A small evaluation reports that one condition beats another. Underneath that
sentence is a table of questions, each marked right or wrong, and the sentence
is true because of how a few of those marks fell. How few is a question with an
exact answer, and it is almost never asked.

This paper asks it of one completed evaluation. A question set of \Questions{}
items was answered by an instruction model under \ConditionsWord{} retrieval
conditions and graded against gold strings, twice over --- once on widely
documented entities and once on less documented ones --- for \Answers{} graded
answers. The project that produced those runs computed \Contrasts{} paired
contrasts from them and recorded which it considered to have found something.
Nothing here re-runs any of that. The rows are on disk, the recorded contrasts
are on disk, and the only new object is a count: for each conclusion, the
smallest number of individual gradings that would have to have come out the
other way for it to stop holding.

The count is exact rather than estimated. A paired contrast between two
conditions reduces to four numbers --- how many questions both got right, how
many each got right alone, how many neither got right --- and every quantity
computed from it, including the exact test, is a function of those four. One
grading change moves one question from one of the four to an adjacent one. So
the tables reachable by $k$ grading changes are exactly those reachable in $k$
steps, the space is small, and a breadth-first walk returns the true minimum.

Three conclusions were named before the search was written, and the answers are
not alike. The reported finding that the multimodal condition trails the
text-only one on image-dependent questions rests on \HarmK{} grading, and
\HarmCandidates{} of the \HarmGradings{} individual gradings on that subset
would each on their own erase it. The reported \HelpDiff{}-point difference in
the other direction needs \HelpK{} gradings to reverse. The one contrast that
survives correction for multiplicity needs \SigK{}. The prespecified kill for
this study was that all three come out at \KillThreshold{} or more, in which
case the claim that these conclusions turn on individual records would have
failed; it \KillState{}, and the smallest is \SmallestK{}.

Two things make this more than an arithmetic exercise. The first is that the
single grading is not hypothetical. An independent regrade of the same answers
by three graders finds exactly \RegradeContested{} answer they disagree about,
and it is one of the \HarmCandidates{}: applied as the record states it, the
reported difference moves from \HarmDiff{} to \RegradedDiff{} percentage points.
The conclusion and its negation are separated by one adjudication of one string.

The second is that a contrast carries two conclusions, not one, and they are not
equally secure. The \HelpDiff{}-point difference needs \HelpK{} gradings to
change direction, which is comparatively sturdy at this scale. Its verdict ---
that an exact test does not call it significant --- needs \HelpVerdictK{}. On
the same four counts, from the same \SubsetN{} questions, the direction is
\HelpRatio{} times further from being overturned than the significance statement
attached to it. Among the \Directional{} contrasts that have a direction to
reverse at all, the two distances differ by \DecoupledGap{} gradings or more in
\DecoupledCount{} of them. Reporting a contrast as ``not significant''
and reporting it as ``smaller'' are not the same claim and are not equally
supported, and nothing in a conventional results table distinguishes them.

A third observation falls out of the same rows and is arguably the most
uncomfortable. Of the \Contrasts{} contrasts, \Incapable{} have a smallest
attainable p-value above the threshold: the two conditions disagreed on so few
questions that no assignment of those disagreements could have produced a
significant result. Their non-significance is a fact about the question set and
carries no information about the conditions being compared. A further
\AtFloor{} returned exactly the most extreme value available to them, so they
were at the limit of what the design could say.

None of this adjudicates the underlying comparison, and this paper makes no
claim about whether the retrieval conditions differ. The reference benchmark
this line of work would be measured against was never obtained, the tier that
would hold that comparison is recorded as \PrimaryTier{} and the tier above it
as \SotaTier{}, and the counting below does not stand in for either. The
contribution is the count and what it exposes: at this scale a conclusion is a
property of individual gradings, and a results table that prints a difference
and a p-value while omitting how many gradings each rests on has withheld the
part a reader would need in order to know what they are looking at.

\section{Methods}
%% Independent methods file. Not woven into main.tex prose.
%% Every quantity described here resolves to a hash-bound entry in the manifest
%% before it may be cited in Results.

\subsection{What was already on disk}

Nothing in this study was generated for it. No model was run, no answer was
produced, and no question was written. Every quantity below comes from records
written by runs that finished before this analysis began, and each record is
bound to the manuscript by its SHA-256 digest; the analysis code re-hashes each
file before reading it and stops the build on any mismatch.

The evaluation being re-read is small and simple. A question set of
\Questions{} questions was answered by a \Model{} instruction model at
temperature \Temperature{} under \ConditionsWord{} retrieval conditions --- no
retrieval, retrieval over a text-only knowledge graph, and retrieval over the
same graph with image-derived attributes attached --- and each answer was
graded right or wrong against a gold string. The latter two are the
graph-structured form of retrieval-augmented
generation~\cite{lewis2020retrieval,edge2024local} and differ only in whether
the graph carries attributes derived from images. Two such question sets exist, one
built around widely documented entities and one around less documented ones,
giving \Answers{} graded answers in total. Each set is split by a recorded
per-question flag into the \SubsetN{} questions whose answer depends on the
image and the \TextOnlyN{} that can be answered from text alone.

The design is paired: the same question is answered under every condition, so a
contrast between two conditions is decided entirely by the questions the two
conditions disagree on. Everything in this paper follows from that.

\subsection{The four counts a contrast reduces to}

For two conditions over a set of questions, each question falls into exactly one
of four cells: both conditions right, the first right alone, the second right
alone, or neither right. The difference in accuracy, the exact test and its
interval are all functions of those four counts and of nothing else.

Changing one grading --- deciding that one answer under one condition should
have been marked the other way --- moves one question from one cell to an
adjacent one. So the set of tables reachable from the observed one by $k$
grading changes is exactly the set reachable in $k$ steps of that move, and the
question this paper asks has an exact answer rather than an estimate.

\subsection{Two conclusions per contrast, and two distances}

A contrast can carry two different claims, and they are unseated by different
things. The first is a direction: which condition is ahead. The second is a
verdict: whether an exact test calls the difference significant at the
conventional threshold of \AlphaLevel{}. This study measures the distance to
each separately, because its central observation is that they are frequently not
the same distance.

The direction is unseated when the difference reaches zero or reverses. The
verdict is unseated when the p-value crosses the threshold, in whichever
direction it currently sits: a significant result is unseated by becoming
non-significant, and a non-significant one by becoming significant. The second
half of that is not symmetry for its own sake. A verdict of ``no effect'' is
reported and acted on exactly as a verdict of ``effect'' is, so it is entitled
to the same audit.

\subsection{Finding the minimum, rather than bounding it}

The search is exhaustive over reachable states rather than greedy over answers.
Because a contrast is fully described by four non-negative counts summing to
$n$, the reachable tables form a small graph, and a breadth-first walk outward
from the observed table returns the first state satisfying the stopping
condition. That state is at the true minimum distance, so the numbers reported
here are minima and not upper bounds on minima.

Where the minimum is one, the paper enumerates every individual grading that
would suffice on its own, so a reader can see which answers carry a conclusion
rather than being told how many. Several candidates do not make a conclusion
sturdier: it means there are several answers, any one of which carries it alone.

\subsection{What the test could have returned}

An exact test on paired binary data is a binomial test at one half on the
discordant pairs alone~\cite{mcnemar1947note}, which is the standard choice for
comparing two systems on a shared test set~\cite{dror2018hitchhiker}, so the
smallest p-value it can
return is fixed by the number of discordant pairs before any answer is graded.
That floor is reported for every contrast. A contrast whose floor lies above the
threshold could not have been called significant however the answers came out,
and its non-significance is therefore a fact about the question set rather than
about the conditions being compared.

Intervals are Clopper--Pearson on the discordant split mapped onto the paired
difference~\cite{clopper1934use}, matching the construction the records already
carry. Where a contrast has no discordant pairs the test is undefined and its
p-value is reported as one rather than as an error.

\subsection{Deletion, as a separate question}

Beside each flip distance the paper reports a leave-one-out
count~\cite{efron1979bootstrap}: how many of the $n$ single-question deletions
leave the direction, and separately the verdict, as it was. This measures
something different. A flip asks what a different grading of the same question
set would have done; a deletion asks what a different question set would have
done. A conclusion can be robust to one and fragile to the other, and both are
reported for every contrast rather than one standing in for the other.

\subsection{Multiplicity}

The records report \Contrasts{} paired contrasts and state that their p-values
are unadjusted. This paper adds the adjustment the record names as the one that
would raise the bar~\cite{holm1979simple} and reports it beside each raw value.
The adjustment is applied to the whole family of \Contrasts{}, because that is
the family that was computed; no subset of it is treated as though it had been
specified in advance.

\subsection{Checks the build enforces}

Several sentences in this paper are enforced as build-time checks rather than
asserted in prose, and the analysis code stops rather than reporting a number
that would contradict them.

Every one of the \Contrasts{} contrasts is recomputed from the per-answer rows
and checked against the recorded version on three quantities separately: the
number of questions, the difference in accuracy, and the exact p-value. The
discordant split is checked as well as the difference, because a table matching
on the difference alone could produce the same p-value by coincidence. The
recorded verdict is checked to follow from the recorded test rather than being
taken on its word.

The reshaping from per-answer rows into paired tables refuses a question that
was not answered under every condition, refuses a condition that answers the
same question twice, and refuses a question whose subset flag differs between
conditions --- each of which would let a pair be formed across the wrong things.
Each run's own summary is checked to still restate the rows in its own
directory, and the fixed measurement table the project recorded as its result is
checked to still come from those rows.

The regrade record is checked to contain exactly one contested grading, and to
be the one the manuscript names. The value the project recorded as not to be
used is checked to still be recorded that way. The tiers reported here as absent
or incomplete are checked to still be in that state, the blocked component is
checked to still be blocked with no score substituted for it, and the retrieval
result is checked to still carry the index state that keeps it a footnote.

\subsection{What is deliberately not computed}

No answer is regenerated and no grading is decided here; the contested grading
discussed below is one the records already contain, applied as the record states
it, not a judgement made by this paper. The reference benchmark this line of
work would eventually be measured against was never downloaded, the tier that
would hold that comparison is recorded as \PrimaryTier{} and the tier above it
as \SotaTier{}, and no value is imputed for either. The \BlockedComponent{}
component is recorded as blocked on a service that was unreachable, and its
scores \ForgedState{} invented in its place. Nothing in this paper substitutes
for any of them.


\section{How far each conclusion is from being overturned}

Table~\ref{tab:contrasts} gives all \Contrasts{} contrasts and
Figure~\ref{fig:flips} plots the two distances for each. Every recorded
difference and every recorded p-value reproduced from the per-answer rows, and
so did every discordant split; the build would have stopped otherwise, and what
follows is therefore a re-reading of the project's own numbers rather than a
second measurement of them.

The three prespecified conclusions are in Table~\ref{tab:prereg}. Taking them in
turn.

The first is that on the image-dependent questions of the widely-documented set,
the multimodal condition is behind the text-only one: \HarmMM{} against
\HarmText{} over \SubsetN{} questions, a difference of \HarmDiff{} percentage
points on \HarmDiscordant{} discordant pairs, with an exact p-value of
\HarmP{}. The minimum flip set is \HarmK{} grading, and the enumeration is
worth more than the count: \HarmCandidates{} of the \HarmGradings{} individual
gradings on that subset would each, alone, be enough to remove the difference or
reverse it. Slightly more than half of the individual judgements underlying that
subset are each independently sufficient to erase its reported direction.

The second is that on the image-dependent questions of the less-documented set,
the multimodal condition is ahead: \HelpMM{} against \HelpText{}, a difference
of \HelpDiff{} percentage points, exact p-value \HelpP{}, which the records
report as not significant and which this paper does not treat as an established
effect. Its direction needs \HelpK{} gradings to reverse. That is the sturdiest
direction among the three at this subset size, and it is still \HelpK{} answers
out of \SubsetNWord{}.

The third is the only contrast in the study that survives adjustment across the
family of \Contrasts{}: on the text-answerable questions of the less-documented
set, the text-retrieval condition beats no retrieval by \SigDiff{} percentage
points with an exact p-value of \SigP{}, or \SigHolm{} after correction. Its
verdict needs \SigK{} gradings and its direction needs \SigSignK{}. This is what
a comparatively secure result looks like at this scale, and it is also the least
interesting one available: it says that supplying a fact makes a model more
likely to state that fact.

\begin{table}[tbp]
\centering
\caption{\CapPreregTable}
\label{tab:prereg}
\small
% Generated by the figure script from hash-bound evidence. Do not hand-edit.
\begin{tabular}{lrrrr}
\toprule
 & As & Gradings & Single-grading & Deletions \\
Conclusion & measured & to unseat & candidates & survived \\
\midrule
Multimodal is behind text, famous & $-5.6$ & 1 & 19 & 14 of 18 \\
Multimodal is ahead of text, long-tail & $+33.3$ & 6 & --- & 18 of 18 \\
Text beats no retrieval, long-tail & $+60.0$, $4.9 \times 10^{-4}$ & 4 & --- & 20 of 20 \\
\midrule
Same contrast, verdict not direction & 0.070 & 1 & 12 & 17 of 18 \\
\bottomrule
\end{tabular}

\end{table}

\begin{figure}[tbp]
\centering
\includegraphics[width=\linewidth]{fig1_flips.pdf}
\caption{\CapFlips}
\label{fig:flips}
\end{figure}

\subsection{The direction and the verdict are not equally secure}

The fourth row of Table~\ref{tab:prereg} is not one of the three, and it is the
result this paper would keep if it could keep only one. The
\HelpDiff{}-point difference whose direction needs \HelpK{} gradings has a
verdict --- not significant at \AlphaLevel{} --- that needs \HelpVerdictK{}, and
\HelpVerdictCandidates{} distinct single gradings would each be enough to make
it significant. Its floor is \HelpFloor{}, so it was capable of significance and
simply did not reach it.

The same four counts, read two ways, give a claim that \HelpK{} answers could
not shake and a claim that \HelpVerdictK{} answer could. Which of the two a
reader takes away is decided by the conventions of the table it is printed in,
and a table reporting $p = \HelpP{}$ with no further annotation invites the
reading that has the weaker support.

Across the study this is the common case rather than the exception. In
\DecoupledCount{} of the \Directional{} contrasts that have a direction to
reverse, the two distances differ by \DecoupledGap{} gradings or more, and they
differ in both directions: \FragileSign{} contrasts have a direction that
\FragileMax{} gradings or fewer would reverse, \FragileVerdict{} have a verdict
that \FragileMax{} or fewer would flip, and no contrast is in both groups.
There is no general rule that one is sturdier. There is only the
fact that they are different questions, and that reporting one number for both
loses the distinction.

\subsection{The single flip is on record}

Figure~\ref{fig:regrade} makes the first conclusion concrete. The same
\RegradeAnswers{} answers were scored by \Graders{} graders, which agree on all
but \RegradeContested{}. The disagreement is question \RegradeQuestion{} under the
multimodal condition: the model answered \RegradeResponse{} where the gold value
is \RegradeGold{}, which a strict string match marks wrong and a
numeric-tolerant match marks right. It is one of the \HarmCandidates{} gradings
the search identified.

Applying it as the record states moves the multimodal condition from
\HarmMM{} to \RegradedMM{} on those \SubsetN{} questions and the contrast from
\HarmDiff{} to \RegradedDiff{} percentage points. The reported direction does
not weaken; it disappears. Under one grader the multimodal condition is behind,
under the other the two are exactly level, and the two graders differ on how to
read one numeric string.

This paper does not endorse either grader, and the choice between them is a real
methodological question rather than a formality --- whether a tolerance is
appropriate depends on what the question was asking, and that is a judgement
about the item. The point is upstream of the choice. A reported direction that
is decided by an unstated grading convention was never a property of the
conditions being compared, and the way to notice that before publishing is to
count, because the count is what makes a grading convention visible as a
load-bearing part of the result.

That the flip turned up in an independent regrade rather than in the search is
also the reason to trust the search. The two procedures had nothing to do with
each other: one enumerated gradings that would change a conclusion, the other
re-scored answers under different string-matching rules, and they met at the
same answer.

\begin{figure}[tbp]
\centering
\includegraphics[width=\linewidth]{fig2_regrade.pdf}
\caption{\CapRegrade}
\label{fig:regrade}
\end{figure}

\section{What the tests could have returned}

Figure~\ref{fig:floor} reports a quantity that is available before any answer is
graded. An exact paired test is a binomial test on the discordant pairs
alone~\cite{mcnemar1947note}, so the smallest p-value it can produce is fixed by
how many questions the two conditions disagreed on. With three discordant pairs
the most extreme result attainable is a two-sided p-value of one quarter, and no
grading of those three answers can do better.

Of the \Contrasts{} contrasts, \Incapable{} have a floor above \AlphaLevel{}.
They were reported as non-significant, and they could not have been reported as
anything else. For those contrasts the phrase carries no information about the
conditions being compared: it is a restatement of how rarely the two conditions
disagreed, which is a property of the question set.

A second reading of the same figure is the \AtFloor{} contrasts sitting on the
diagonal, where every discordant pair fell the same way and the test returned
the most extreme value available to it. Of those, \AtFloorNS{} are
non-significant regardless. A contrast in which one condition won every single
question the two conditions disagreed about, and which is nonetheless reported
as no evidence of a difference, is the clearest available demonstration that a
threshold verdict at this scale is a statement about $n$ before it is a
statement about the conditions.

Of the \Contrasts{} contrasts, \SignificantRaw{} reach the threshold unadjusted
and \SurvivesHolm{} survive correction across the family~\cite{holm1979simple}.
The records state that their p-values are unadjusted and that a correction would
raise the bar; this applies it. The survivors are all differences against no
retrieval at all, and none of them is a comparison between the two retrieval
conditions, which is the comparison the evaluation was built to make.

\begin{figure}[tbp]
\centering
\includegraphics[width=0.62\linewidth]{fig3_floor.pdf}
\caption{\CapFloor}
\label{fig:floor}
\end{figure}

\begin{table}[tbp]
\centering
\caption{\CapContrastsTable}
\label{tab:contrasts}
\scriptsize
\setlength{\tabcolsep}{4.5pt}
% Generated by the figure script from hash-bound evidence. Do not hand-edit.
\begin{tabular}{llrrrrrrrr}
\toprule
& & & & & & & \multicolumn{2}{c}{Gradings to unseat} & \\
\cmidrule(lr){8-9}
Questions & Contrast & $n$ & Disc. & Diff. (pt) & $p$ & Floor & Dir. & Verd. & Holm \\
\midrule
Famous, image & multimodal vs text & 18 & 7 & $-5.6$ & 1.000 & 0.016 & 1 & 5 & 1.000 \\
Famous, image & text vs none & 18 & 2 & $+0.0$ & 1.000 & 0.500 & 0 & 6 & 1.000 \\
Famous, image & multimodal vs none & 18 & 7 & $-5.6$ & 1.000 & 0.016 & 1 & 5 & 1.000 \\
Famous, text & multimodal vs text & 20 & 0 & $+0.0$ & 1.000 & 1.000 & 0 & 6 & 1.000 \\
Famous, text & text vs none & 20 & 4 & $+20.0$ & 0.125 & 0.125 & 4 & 2 & 1.000 \\
Famous, text & multimodal vs none & 20 & 4 & $+20.0$ & 0.125 & 0.125 & 4 & 2 & 1.000 \\
Famous, all & multimodal vs text & 38 & 7 & $-2.6$ & 1.000 & 0.016 & 1 & 5 & 1.000 \\
Famous, all & text vs none & 38 & 6 & $+10.5$ & 0.219 & 0.031 & 4 & 2 & 1.000 \\
Famous, all & multimodal vs none & 38 & 11 & $+7.9$ & 0.549 & $9.8 \times 10^{-4}$ & 3 & 4 & 1.000 \\
Long-tail, image & multimodal vs text & 18 & 8 & $+33.3$ & 0.070 & 0.008 & 6 & 1 & 0.984 \\
Long-tail, image & text vs none & 18 & 4 & $-11.1$ & 0.625 & 0.125 & 2 & 4 & 1.000 \\
Long-tail, image & multimodal vs none & 18 & 8 & $+22.2$ & 0.289 & 0.008 & 4 & 2 & 1.000 \\
Long-tail, text & multimodal vs text & 20 & 0 & $+0.0$ & 1.000 & 1.000 & 0 & 6 & 1.000 \\
Long-tail, text & text vs none & 20 & 12 & $+60.0$ & $4.9 \times 10^{-4}$ & $4.9 \times 10^{-4}$ & 12 & 4 & 0.008 \\
Long-tail, text & multimodal vs none & 20 & 12 & $+60.0$ & $4.9 \times 10^{-4}$ & $4.9 \times 10^{-4}$ & 12 & 4 & 0.008 \\
Long-tail, all & multimodal vs text & 38 & 8 & $+15.8$ & 0.070 & 0.008 & 6 & 1 & 0.984 \\
Long-tail, all & text vs none & 38 & 16 & $+26.3$ & 0.021 & $3.1 \times 10^{-5}$ & 10 & 2 & 0.319 \\
Long-tail, all & multimodal vs none & 38 & 20 & $+42.1$ & $4.0 \times 10^{-4}$ & $1.9 \times 10^{-6}$ & 16 & 6 & 0.007 \\
\bottomrule
\end{tabular}

\end{table}

\section{Deleting a question is not re-grading an answer}

Figure~\ref{fig:loo} reports the leave-one-out count beside the flip distance,
and the two disagree in an instructive way. Most contrasts survive every single
deletion: the direction and the verdict are unchanged in all $n$ recomputations.
Deletion is a blunt instrument here, because removing a concordant question
changes neither the discordant split nor the sign, and most questions are
concordant.

The exceptions are the contrasts already identified. The direction of the first
conclusion survives \HarmLooSign{} of \SubsetN{} deletions, so \HarmLooBreaks{}
individual questions each carry it by their presence as well as by their
grading. The verdict on the second survives \HelpLooVerdict{} of \SubsetN{}.

The two measurements answer different counterfactuals and should be read as
such. A flip asks what a different adjudication of the same evaluation would
have concluded, which is a question about grading policy and about annotator
disagreement~\cite{northcutt2021pervasive}. A deletion asks what a slightly
different question set would have concluded, which is a question about sampling.
An evaluation can be exposed on one and not the other, and reporting only the
one that happens to look better is a degree of freedom worth closing in
advance~\cite{simmons2011false}.

\begin{figure}[tbp]
\centering
\includegraphics[width=\linewidth]{fig4_loo.pdf}
\caption{\CapLeaveOneOut}
\label{fig:loo}
\end{figure}

\section{Discussion}

\paragraph{The count is old; the target is not.} Counting how many outcomes
would have to change to overturn a result is the fragility
index~\cite{feinstein1990unit,walsh2014statistical}, introduced for randomised
trials, where it revealed that a large share of results in the clinical
literature turn on a handful of events --- one mechanism behind the more general
observation that small studies produce unreliable
findings~\cite{ioannidis2005why}. The construction transfers to machine
learning evaluation almost unchanged, and the reasons to want it are stronger
here: an evaluation item is graded by a rule that someone wrote, benchmark labels
are known to carry errors at rates that change published
rankings~\cite{northcutt2021pervasive}, and the evaluation sets that appear in
practice are frequently small enough that the count is in single
digits~\cite{card2020little}. The main criticism of the index in its original
setting is that it is a monotone function of the p-value and so adds
nothing~\cite{carter2017fragility}. That criticism has less force here, for a
reason this study demonstrates rather than argues: the count is being applied to
the direction of a difference as well as to the verdict on it, the two are not
monotone in each other, and \DecoupledCount{} of the \Directional{} contrasts
that have a direction to reverse separate them by \DecoupledGap{} gradings or
more.

\paragraph{Two conclusions per contrast, and only one of them gets audited.}
Significance is routinely questioned and directions are routinely believed. This
study found the opposite pattern at least as often. The clearest case is the
contrast whose direction \HelpK{} gradings could not reverse and whose verdict
\HelpVerdictK{} grading would flip, but the general point is that these are separate claims with
separate supports, and a results table reports them in a single row as though
they stood or fell together. The recommendation implied is small: print the two
counts. They cost nothing to compute, they are exact, and they say what a
confidence interval says while remaining in the units the reader can inspect ---
answers.

\paragraph{Non-significance is not a finding when significance was
unreachable.} Of the \Contrasts{} contrasts here, \Incapable{} could not have
been significant at any grading of their answers. Reporting them as null results
is not conservative; it attributes to the conditions a property of the question
set. The floor is one line of arithmetic from the discordant count and belongs
beside any negative result from a paired design, in the same way that a power
calculation belongs beside a negative result from an unpaired
one~\cite{card2020little,amrhein2019scientists}. The stronger version of the
point is the contrast in which every disagreement fell one way and the verdict
was still negative: unanimity among the informative questions was not enough,
because there were \AtFloorNSDiscordant{} of them.

\paragraph{A grading convention is a researcher degree of freedom.} String
matching against gold answers looks like a mechanical step and is not one.
Strict matching, substring matching and numeric tolerance are three defensible
rules, they disagree about one answer here, and that one answer decides a
reported direction. The literature on undisclosed analytic
flexibility~\cite{simmons2011false,gelman2014statistical} is usually read as
being about model selection and outlier handling; grading rules belong on the
same list, and in small evaluations they may dominate it. The regrade in this
study is the practice that made the point visible, and it cost one pass over the
stored answers with no model in the loop.

\paragraph{Report the count with the claim.} The concrete proposal is one column.
Beside each reported contrast, print the smallest number of gradings that would
overturn its direction and the smallest that would move its verdict. Both are
exhaustively computable at these sizes. A reader seeing \HarmK{} in that column
knows immediately what kind of object they are being shown, without needing to
reconstruct the discordant table from an accuracy and a sample size --- which is
what it currently takes, and which is why the question is rarely asked. This is
the same instinct behind reporting effect sizes and intervals rather than
verdicts~\cite{wasserstein2016asa,amrhein2019scientists}, expressed in the unit
that a small evaluation is actually made of.

\paragraph{What this does not fix.} Counting gradings makes fragility visible
and does nothing about its cause. An evaluation whose conclusions rest on single
answers needs more questions, not better reporting of how few it
had~\cite{card2020little,bowman2021will}. The count also says nothing about
whether the question set measures what it claims to measure, which at these
sizes is the larger risk~\cite{raji2021everything}. Its value is narrower: it is
cheap, it is exact, it applies to records that already exist, and it prevents
one specific error, which is publishing an arithmetic accident in the register
of a finding.

\section{What this study does not claim}

No claim is made that adding image-derived attributes to retrieval helps or
harms. The contrast on the less-documented set is \HelpDiff{} percentage points
with an exact p-value of \HelpP{}, which does not reach the threshold and is not
treated here as an established effect; the contrast on the widely-documented set
is \HarmDiff{} points and is overturned by \HarmK{} grading. Both are reported
as objects whose fragility is being measured, and neither is offered as
evidence about the underlying conditions.

No claim is made about performance on any reference benchmark. The benchmark
this line of work would be measured against was never obtained, the tier that
would hold that comparison is recorded as \PrimaryTier{} and the tier above it
as \SotaTier{}, and no value is imputed for either. No result from any other
system is introduced for comparison.

No retrieval result here is a benchmark accuracy. The retrieval step resolves
\RetrievalCorrect{} of \RetrievalN{} questions to the right entity on each
question set, and that figure is recorded as a footnote on an easy
small-graph lookup with \IndexFiles{} files in the index it names. It is
reported here only to establish that the conditions differ in what they
retrieved rather than in whether retrieval worked at all.

No score is supplied for the blocked component. The \BlockedComponent{}
component is recorded as blocked on a service that was unreachable, its scores
\ForgedState{} invented in its place, and nothing here substitutes for them. The
tier that would hold the full comparison remains \PrimaryTier{} and this paper
does not present its own counting as progress toward completing it.

No use is made of the value the project recorded as not to be used. An earlier
run of the same subset under a different caption source produced a much higher
figure for the multimodal condition; the records mark it as unused and name the
real-caption values as the result, and this paper quotes the real-caption values
only. The excluded figure appears nowhere in this manuscript.

No grading is adjudicated here. The contested answer is contested in the records
already, both readings of it are the records' own, and this paper reports what
each implies rather than deciding between them. Nothing was re-graded, re-run or
re-scored for this study.

No claim is made that a minimum flip set of one means a conclusion is false. It
means the conclusion is supported by one grading. A true effect measured on
eighteen questions can also rest on one grading, and the count does not
distinguish those cases; it only establishes what would be needed to distinguish
them, which is more questions.

\section{Conclusion}

One completed evaluation of \Questions{} questions under \ConditionsWord{} retrieval
conditions was re-read by counting rather than by re-running. For each of its
\Contrasts{} paired contrasts, the smallest number of individual gradings that
would overturn the direction, and separately the verdict, was found by
exhaustive search over the reachable tables.

The three conclusions named in advance came out at \HarmK{}, \HelpK{} and
\SigK{} gradings, so the prespecified kill \KillState{}. The smallest,
\SmallestK{}, belongs to the reported finding that the multimodal condition
trails the text-only one: \HarmCandidates{} of the \HarmGradings{} gradings on
that subset would each erase it alone, and an independent regrade of the same
answers disagrees about exactly one of them, which moves the difference from
\HarmDiff{} to \RegradedDiff{} percentage points. The direction of the opposing
\HelpDiff{}-point difference needs \HelpK{} gradings while the verdict attached
to it needs \HelpVerdictK{}, and across the study those two distances differ by
\DecoupledGap{} or more in \DecoupledCount{} of the \Directional{} contrasts
that have a direction at all. No grading of the answers could have made
\Incapable{} of the contrasts significant.

The recommendation is one column in a results table: beside each contrast, the
number of gradings that would overturn its direction and the number that would
move its verdict. It is exactly computable, it costs nothing at these sizes, and
it converts a question a reader cannot currently ask into a number they can
read. When the count is large, nothing has been lost by printing it. When it is
\HarmK{}, what has been learned is that the row is reporting an adjudication of
one answer, and that is worth knowing before the table is published rather than
after.

\section*{Data and code availability}

The per-answer records for both question sets, the summaries each run wrote for
itself, the recorded paired contrasts with their exact tests and intervals, the
independent regrade under three graders, the fixed measurement table, the
recorded state of each comparison tier including those that are absent or
blocked, and the analysis code that produces every number in this manuscript
accompany the manuscript source. Each artifact is recorded with its SHA-256
digest in a manifest, and the analysis code verifies those digests against the
bytes on disk before reading them, so the build fails rather than silently
reporting stale numbers. The knowledge graph and question sets are small
local artifacts distributed with the records; no third-party benchmark is
redistributed, and the model used to produce the answers is a third-party
release identified by name rather than copied.
%% Generated by scripts/release_targets.py from the release ledger.
%% Do not edit: an identifier typed here would not be checked against
%% anything, which is exactly the failure the ledger exists to prevent.
%% Ledger state: github=CREATED, zenodo=DEPOSITED
That material is deposited under the persistent identifier \href{https://doi.org/10.5281/zenodo.22647028}{10.5281/zenodo.22647028}, and the development repository is at \url{https://github.com/PeterPonyu/minimum-flip-set-audit}.


\section*{Ethics statement}

This study is a secondary analysis of numeric records produced by earlier
computational runs. No human subjects are involved, no personal data are
processed, and no new model output was generated. The questions concern publicly
documented landmarks and comparable entities. One consideration bears directly
on the subject matter: the paper's central observation is that a reported
difference between two systems can be decided by a single grading decision, and
the same is true of evaluations used to support claims about deployed systems,
where a fragile result may be read as a safety or capability assurance.
Reporting the count is offered partly for that reason.

\section*{Competing interests}

Not yet declared. This is a working draft and the declaration is outstanding.

\section*{Funding and author contributions}

Not yet declared. This is a working draft; the author list is not fixed and
contributions are therefore not yet attributable.

\section*{Declaration of generative AI use}

The analysis code, the figure code and the drafting of this manuscript were
produced with substantial assistance from large language model agents working
under human direction. Two safeguards bear directly on what is reported. Every
number printed above is emitted programmatically from artifacts bound by
cryptographic digest rather than transcribed by a model into prose. And several
of the paper's structural claims --- that every recorded difference, discordant
split and exact p-value reproduces from the per-answer rows, that each run's
summary still restates the rows in its own directory, that the regrade record
contains exactly one contested grading and that it is the one named here, and
that the tiers reported as absent or blocked are still in that state with no
score substituted --- are enforced as build-time checks that stop the build
rather than as sentences. The answers being analysed were themselves produced by
a language model in the earlier runs this study re-reads; none was regenerated
here. Responsibility for the content rests with the authors.

\bibliographystyle{plain}
\bibliography{references}

\end{document}
